% THE Q-MODEL — Final Research Monograph
% Version 1.0 (final) — Publication-quality scientific monograph
% Build: pdflatex TQM_v1_0_Monograph_Final.tex (twice)
%
% This is the final editorial pass: higher information density, compact layout,
% boxed summaries, per-theorem historical context, per-equation reading notes,
% a framework-evolution chapter, and a hostile-review survival table.
% It adds NO new physics, NO new derivations, and changes NO conclusions.
% Publication status (updated): READY AS A FOUNDATION MONOGRAPH —
% READY AS A RESEARCH-PROGRAM PAPER — NOT READY AS A DERIVATION PAPER.

\documentclass[10pt,twoside,openany]{book}
\usepackage[utf8]{inputenc}
\usepackage[T1]{fontenc}
\usepackage{amsmath,amssymb,amsthm}
\usepackage[margin=2.4cm]{geometry}
\usepackage{booktabs}
\usepackage{array}
\usepackage{longtable}
\usepackage{lmodern}
\usepackage[hidelinks]{hyperref}
\usepackage{xcolor}

\definecolor{caveat}{HTML}{7A2E2E}

% Theorem-like environments
\newtheorem{theorem}{Theorem}[chapter]
\newtheorem{proposition}[theorem]{Proposition}
\newtheorem{lemma}[theorem]{Lemma}
\newtheorem{corollary}[theorem]{Corollary}
\newtheorem{definition}[theorem]{Definition}
\newtheorem{axiom}[theorem]{Axiom}
\newtheorem{remark}[theorem]{Remark}
\newtheorem{nogo}[theorem]{No-Go Theorem}

\newcommand{\derived}{\textbf{DERIVED}}
\newcommand{\realunder}{\textbf{REAL-UNDERIVED}}
\newcommand{\drawn}{\textbf{DRAWN}}

% Reader notes appended to major results: intuition, significance, limitations
\newcommand{\theoremnotes}[3]{%
  \par\medskip\noindent
  \textbf{Intuition.}\ #1\par
  \textbf{Significance.}\ #2\par
  \textbf{Limitations.}\ #3\par}

% Compact part divider (inline heading, no full-page break; keeps numbering)
\makeatletter
\renewcommand{\part}[1]{%
  \refstepcounter{part}%
  \par\addvspace{1.0em}\noindent
  {\centering\Large\bfseries \partname\ \thepart\ ---\ #1\par}%
  \addvspace{0.5em}%
  \addcontentsline{toc}{part}{\protect\numberline{\thepart}#1}}
\makeatother

% mainmatter without a forced right-hand page (removes the blank page before Chapter 1)
\makeatletter
\renewcommand{\mainmatter}{%
  \if@openright\cleardoublepage\else\clearpage\fi
  \@mainmattertrue
  \pagenumbering{arabic}}
\makeatother

% Compact unnumbered part divider
\newcommand{\compactpart}[1]{%
  \par\addvspace{1.0em}\noindent
  {\centering\Large\bfseries #1\par}%
  \addvspace{0.5em}%
  \addcontentsline{toc}{part}{#1}}

% Boxed summary
\newcommand{\summarybox}[1]{%
  \par\medskip\noindent
  \fcolorbox{black}{gray!10}{%
    \begin{minipage}{0.955\textwidth}\small #1\end{minipage}}}

% Chapter-end summary (Key results / Open questions / Verification status)
\newcommand{\chapterend}[3]{%
  \summarybox{%
    \textbf{Key results.}\ #1\par
    \textbf{Open questions.}\ #2\par
    \textbf{Verification status.}\ #3}}

% Equation reading notes (Intuition / Meaning / Why it matters / Limitations)
\newcommand{\eqnotes}[4]{%
  \par\smallskip\noindent
  \textbf{Intuition.}\ #1\ 
  \textbf{Meaning.}\ #2\ 
  \textbf{Why it matters.}\ #3\ 
  \textbf{Limitations.}\ #4\par}

% Theorem historical notes
\newcommand{\historynotes}[4]{%
  \par\smallskip\noindent
  \textbf{Historical context.}\ #1\par
  \textbf{Why it mattered.}\ #2\par
  \textbf{What was tried before.}\ #3\par
  \textbf{Why the result was accepted.}\ #4\par}

\begin{document}

\hypersetup{pdftitle={THE Q-MODEL --- From Q to Cosmology},
pdfauthor={Fabrice Wieser}}

\frontmatter

% ---- Title block (compact; title, author, date on one page, no page break) ----
\begin{center}
{\LARGE\bfseries THE Q-MODEL — From Q to Cosmology\par}
\vspace{0.4em}
{\large A Research Monograph on Structure, Complexity and Random Actualization\par}
\vspace{1.0em}
{\normalsize Fabrice Wieser\par}
{\small Independent Researcher and Full Stack Software Engineer\par}
{\small \url{https://github.com/MagusDraconis/TQM}\par}
\vspace{0.5em}
{\normalsize Version 1.0 (final) \quad·\quad 15 August 2026\par}
\end{center}

\vspace{0.8em}

% ---- Abstract (on the title page, above the publication-status box) ----
\addcontentsline{toc}{chapter}{Abstract}
{\centering\Large\bfseries Abstract\par}
\vspace{0.3em}
\noindent
This monograph is the research-monograph edition of THE Q-MODEL (TQM), a theory of
structure and content that compresses observable physics to four primitives — the
individuation principle $Q$, Random Actualization, the scale triad $(\ell,\tau,\hbar)$,
and a single continuous nonlinearity parameter $M^2$ — and holds that the \emph{form} of
every structure is derivable from these primitives while the \emph{content} is realized by
contingent draw. The monograph presents, in self-contained form: the formal primitives and
dynamical system (the graph Laplacian $L_Q$ and the Schr\"odinger equation); the verified
continuum-limit chain from $L_Q$ to the flat Laplacian, the d'Alembertian, the
Laplace--Beltrami operator, the curved-space Schr\"odinger equation, and the Einstein
tensor; the derivation hierarchy (gauge structure, spatial dimensionality, the
multiplicity bound, the abundance law, phase-gradient gravity); the three-category
taxonomy; the conditional no-go theorems; the quantitative predictions; and the complete
inventory of executable verification results.

In addition to the technical content of the reference edition, this edition adds an
explanatory layer: the historical \emph{origin} of the theory and its transition from its
predecessor framework; a \emph{development timeline} recording the major phases, key
discoveries, and abandoned routes; an account of \emph{why} the structure/content split is
the organizing principle; the \emph{hostile-review history} and what survived it; a
program-by-program \emph{research journey}; and the \emph{lessons learned} from wrong
assumptions, rejected derivations, and no-go discoveries. Throughout, each major theorem is
accompanied by its \emph{intuition}, \emph{significance}, and \emph{limitations}. Nothing
in this narrative layer changes the physics or the conclusions of the reference edition.

\vfill

% ---- Publication status box (pinned to the bottom of the title page) ----
\begin{center}
\fcolorbox{caveat}{white}{%
\begin{minipage}{0.92\textwidth}
\color{caveat}
\textbf{Publication status: READY AS A FOUNDATION MONOGRAPH --- READY AS A
RESEARCH-PROGRAM PAPER --- NOT READY AS A DERIVATION PAPER.}\\[2mm]
A foundation monograph may legitimately use imported proven theorems as inputs; the
remaining blockers affect only a \emph{derivation-paper} claim. This monograph is released
for archival on Zenodo as a foundation monograph and a research-program record; it is
\emph{not} submitted as a peer-reviewed derivation paper. The central derivation claim
(Einstein recovery from Q-events) remains \emph{logical, not mathematical}: the metric and
the BDG action are imported (proven but not TQM-derived), $G=\ell^2c^3/\hbar$ is
dimensional analysis, and no unique sharp prediction yet discriminates TQM from
SM$+\Lambda$CDM. Every claim marked \textbf{verified} is backed by an executable xUnit
test in \texttt{TQM.Tests/ResearchXC} (or \texttt{ResearchQG}); see Part~V and Appendix~A.
\end{minipage}}
\end{center}

\tableofcontents

\mainmatter
\thispagestyle{plain}

%=============================================================================
\compactpart{Genesis and Method}
%=============================================================================

This part records how the theory came to be, why it is organized as it is, and what the
research program actually did. It is historical and methodological, not physical: it
introduces no new claim and revises none. Its purpose is to make the technical parts that
follow legible as the outcome of a specific, auditable research trajectory rather than as
a finished axiomatic system dropped from nowhere.

\chapter{How to Read This Monograph}
\label{ch:reading}

This monograph is a reference edition with an explanatory layer, and the two layers can be
read independently. The \emph{reference} layer (the technical parts and the appendices)
states the theory, its derivations, its classifications, and its verification record. The
\emph{narrative} layer (this part) explains how the theory came to be and why it is
organized as it is. Nothing in the narrative layer changes the physics; it exists to make
the physics legible as the outcome of a specific, auditable program rather than as a
finished axiom system dropped from nowhere.

Three conventions govern the whole document and should be read first.

\textbf{(1) The taxonomy.} Every result is tagged with one of three statuses.
\derived\ means the result is computable from the primitives by a theorem.
\realunder\ (``real-underived'') means the structure is real, precise, and predictive, but
its origin is not computable from the primitives. \drawn\ means the value is a
coincidental draw of the abundance law. The tags are claims \emph{about derivability}, not
about truth: a \realunder\ or \drawn\ result can be exactly right while remaining
underived. Reading a \realunder\ tag as a derivation was the single most common misreading
the hostile review had to correct.

\textbf{(2) What ``verified'' means.} A result marked \textbf{verified} is backed by an
executable xUnit test in the repository. ``PASS'' means the test ran and its assertion
held; it does \emph{not} mean the result was derived from the primitives --- that is the
taxonomy's job. Some tests verify a derivation; others verify a \emph{standard}
mathematical identity (such as the Einstein-tensor chain) precisely in order to show where
the TQM derivation stops.

\textbf{(3) What ``confidence'' means.} The confidence numbers (e.g.\ T-02 $=$ 0.85) are
\emph{uncalibrated ordinal ranks}, not posterior probabilities. They order the program's
own degree of belief; they do not measure it. They are never to be read as $p$-values or
as numerical posteriors.

A short map. This part gives the history and method. The \emph{Foundations} part states
the primitives and the dynamical system. The \emph{Continuum-Limit} part records the two
continuum chains. The \emph{Derivation Hierarchy} part runs dimensionality, gauge, flavor,
gravity, and cosmology. The \emph{Classification} part states the taxonomy, the no-go
theorems, and the predictions. The \emph{Verification} part gives the test record, the
hostile-review audit trail, and the conclusion. The appendices carry the full test
inventory, the confidence tables, a worked derivation, the program ledger, the glossary,
and the references.

\chapter{The Origin of TQM}
\label{ch:origin}

\section{Motivation}

THE Q-MODEL (TQM) began from a single, deliberately simple question: are matter, quantum
behaviour, and gravitation \emph{fundamental}, or do they emerge from something more
basic? The working assumption that has guided the entire program is the claim that
\emph{matter is not fundamental} --- that what we call a particle is a dynamically
stabilized wave structure, and that synchronization and resonance are more fundamental
than particles themselves. From that seed a research path was fixed early and never
reversed:

\begin{center}
Temporal Oscillators $\to$ Synchronization $\to$ Resonance Clusters $\to$
Self-Stabilizing Wave Structures $\to$ Proto-Particles $\to$ Quantum Behaviour $\to$
Gravity $\to$ Cosmology.
\end{center}

Four standing decisions disciplined the program from the outset, and they explain much of
the monograph's structure. \emph{No cosmology first} and \emph{no gravity first} forbade
the temptation to tune a model against the most spectacular data before the microscopic
foundations were in place. \emph{No particle assumptions} forbade importing the Standard
Model's ingredients (gauge groups, generations, masses) as inputs. \emph{Emergence first}
required every ingredient to be earned from below, and \emph{physics before
interpretation} required the dynamical content to be settled before any ontological gloss
was attached. These rules were designed to prevent the theory from being
reverse-engineered toward known physics and to force every element to carry its own
weight.

\section{Predecessor ideas}

TQM did not appear in a vacuum. Its immediate predecessor was the \textbf{Temporal
Resonance Model} (TRM), a framework built on phase-lattice machinery that had already
assembled a number of striking pieces: a \emph{Time Field} (a phase-gradient picture of
gravity), a \emph{Radial Acceleration Relation} (RAR), \emph{Frame Dragging},
a \emph{Memory Channel}, a \emph{Theta Chain}, an \emph{$m=3$ Closure} mechanism (rational
mode-locking $\Omega=(q+3)/q$, $\gamma=1/\Omega$), a \emph{Quantum Engine}, a
\emph{Temporal Drift} (tired-light) module, and a \emph{Unified Action} roadmap. TRM
supplied the conceptual vocabulary --- a temporal field, oscillation, phase, resonance ---
on which the later theory was rebuilt.

Each module deserves a one-line statement of what it actually claimed. The \emph{Time
Field} identified gravity with the phase gradient of a temporal field --- the idea that,
after the transition, became TQM's phase-gradient gravity picture. The \emph{Radial
Acceleration Relation} targeted the galaxy-rotation acceleration scale. \emph{Frame
Dragging} introduced a vector sector (gravitomagnetism). The \emph{Memory Channel} proposed
a history-encoding invariant $\varphi^2|\dot\mu|$. The \emph{Theta Chain} was a nonlocal
observable-extraction chain $\Theta\to O_5\to\lambda_\Theta\to g_{\rm obs}$. The \emph{$m=3$
Closure} proposed rational mode-locking ($\Omega=(q+3)/q$, $\gamma=1/\Omega$) as a mechanism
for the ``why 3'' question. The \emph{Quantum Engine} was a non-unitary quantum module. The
\emph{Temporal Drift} was a tired-light ($\beta_T$) cosmology. The \emph{Unified Action}
was a roadmap $S_{\rm eff}[T,A_T,\Theta]$ depending on all the others. This list matters
because the transition was decided module by module, not wholesale.

The relationship between TRM and TQM is best summarized by the \emph{TRM Legacy Final
Classification}, the systematic audit that determined what, if anything, of the older
framework should migrate into the new one. Of the nine TRM modules, \emph{three} were
\emph{absorbed}: the Time Field became the phase-gradient gravity picture, the RAR became
the zero-parameter relation $g_\dagger = cH_0/2\pi$, and Frame Dragging was recognized as
GR gravitomagnetism (Lense--Thirring, already measured). \emph{Two} were \emph{rejected}:
Temporal Drift (tired light) was falsified by the QG-080--089 reinterpretation audits, and
the Quantum Engine was found to be non-unitary and worse than the BDG lattice operator.
\emph{Three} survived as \emph{candidate mathematics only}: the $m=3$ Closure (a candidate
mechanism for the $N\le3$ upper bound, but unmapped to observables), the Memory Channel
(a genuine invariant $\varphi^2|\dot\mu|$ with no observable), and the Theta Chain (a
homonym of TQM's own information field). \emph{One} remained open: the Unified Action
roadmap.

\section{The TRM $\to$ TQM transition}

The decisive result of the legacy audit was negative and, in its way, liberating: TRM
carried \emph{zero candidate physics} --- no TRM module supplied a genuinely new testable
observable. The residue of value was methodological. What the transition produced was
therefore a \emph{clean reconstruction} rather than an extension: the new theory
re-derived, from the fewest possible ingredients, exactly the structure that TRM had
assembled piecemeal, and discarded the parts that did not survive hostile reduction.

The compressed endpoint of that reconstruction is the subject of this monograph: a theory
with \emph{two irreducible primitives} --- the individuation principle $Q$ and genuine
ontological randomness --- plus \emph{one} contingent continuous number $M^2$ (the
nonlinearity regime), with everything else --- time, space, $3{+}1$ dimensions, Hilbert
space, unitary quantum mechanics, the Schr\"odinger equation, the Born rule, measurement,
gravity as causal structure, particles as topological defects, and $U(1)$ gauge symmetry
--- derived or emergent. The parameter count collapses from the Standard Model's
$\sim19$ free numbers to $2$ primitives plus $1$ number, a reduction of roughly $95\%$,
while $U(1)$ is promoted from an assumption to a theorem.

\chapter{Development Timeline}
\label{ch:timeline}

The program is best read as a sequence of phases, each of which closed one question and
opened the next. This chapter gives the skeleton; the detailed experiment records live in
the project's laboratory book, and the four research programs are narrated individually
in Chapter~\ref{ch:journey}.

\section{Major phases}

The program falls into five successive eras. Each is summarized here in enough detail to
show what was actually done; the four research programs that carry the technical weight
are narrated further in Chapter~\ref{ch:journey}.

\subsection{Early oscillator experiments (TQM-001--TQM-008)}

The first experiments were numerical and exploratory, and their role was to find out
whether the central intuition --- coherent structure emerging from oscillators --- had any
teeth. TQM-001 showed that an order parameter $R$ of a population climbs toward $1$:
coherent synchronization is not put in by hand, it appears. TQM-002 then closed a tempting
blind alley: random matrices give Wigner-like spectra but do \emph{not} produce the kind
of structured emergence the program needed, so randomness alone is not enough. TQM-003
showed that structured topology alone is also insufficient --- the dynamics, not the
static graph, carries the weight. TQM-004 showed that a field-mediated interaction does
not by itself synchronize two asymmetric oscillators; the field behaves more like a
resonance medium than a pure synchronization medium.

TQM-005 found stable resonance clusters of three oscillators forming and persisting over
most of a simulation --- early, tentative support for ``matter as stabilized resonance''
--- but with only a few percent energy participation, so the clusters did not yet dominate
the field. TQM-006 found the first sharp threshold, a critical resonance density
$\rho_c\approx0.09$, above which persistent coherent modes appear: a percolation-like
transition, and the program's first well-defined critical point. TQM-007 reported five
resonance families. TQM-008 then demolished that apparent diversity with a
400-simulation reproducibility study, leaving a single universal coherent mode (F4) and
exposing the others as seed artifacts. The lesson of TQM-008 --- apparent diversity can be
a single seed --- became a standing methodological caution for everything that followed.

\subsection{Alternative foundations (ResearchX, X001--X034)}

A parallel track attacked the foundations from the side, and it is responsible for several
of the theory's deepest results. X001 found eleven hidden assumptions, with the static
graph chief among them. X002--X004 tested dynamism (node motion, mobility, graph growth)
and found no new phenomena: dynamic graphs alone produce no open-ended innovation. X005--
X006 showed that nonlinearity breaks eigenmodes and creates solitons, so the ``quantum''
reading of the early linear results had to be withdrawn --- Q, fitness, and evolution
survived, the naive quantum reading did not. X007--X015 built the general theory of
information carriers and proved that \emph{reversibility} plus \emph{self-consistency} is
the minimally sufficient pair whose intersection is unitary quantum mechanics.

X016--X022 built the complexity staircase and proved that open-ended evolution cannot occur
in a finite system (a pigeonhole bound). X023--X031 then proved that quantum reality is
the unique global maximum of finite complexity --- a necessity result, not an anthropic
one. X032--X034 unified the main line with ResearchX, confirming that $L_Q$ is not
fundamental: any operator satisfying reversibility and self-consistency at their joint
maximum yields unitary quantum mechanics. This is the result behind the claim that the
theory has been compressed to its minimal form.

\subsection{Spacetime emergence (X040--X046)}

Time was derived as the partial order of actualization events ($E_1<E_2$ iff $E_2$ depends
on $E_1$'s outcome), with the arrow of time inherited from the irreversibility of
actualization. Gravity was placed on a causal-set footing, and a reconstruction experiment
recovered distances from Q-event correlations. Dimensionality $3{+}1$ was derived from
complexity maximization (capacity times stability times accuracy). Newton's constant,
$\hbar$, and $\Lambda$ were all traced to the single event scale $\ell$: $G=\ell^2c^3/\hbar$,
$\hbar=1$ in event units, and $\Lambda\sim1/\sqrt N$ from Poisson fluctuations of the
event count. This is the era in which the theory stopped needing to \emph{assume}
spacetime and began to \emph{produce} it.

\subsection{Matter emergence (X047--X060)}

Particles were identified with topological defects (solitons): charge with the Betti
number, mass with defect energy, spin with orientation. $U(1)$ was derived as the
automorphism group of the circle. Generations were characterized as excitation levels of
the same topology, with the mass hierarchy following an anharmonic spacing. Fermion mixing
was traced to defect wavefunction overlap (hierarchical CKM from Dirac defects, anarchic
PMNS from Majorana defects). Neutrino masses were explained by delocalized defects, and
normal ordering was predicted from an attractive self-interaction. The Standard Model
gauge group was shown to be \emph{preferred} by ecological fitness but \emph{not} uniquely
derivable --- the largest remaining gap in the matter program.

\subsection{Parameter compression and closure (Phases 140--159)}

The final era consolidated. The parameter-count audits (X057--X060f) reduced the theory to
two primitives plus one number, with $U(1)$ promoted to a theorem and $M^2$ identified as
the single irreducible continuous parameter. The late phases then closed the residual
questions one by one: the RAR derivation $g_\dagger=cH_0/2\pi$; the Koide closure (seven
routes exhausted, six no-gos plus one blocked); the gauge-origin audit ($U(1)$ derived,
$SU(2)$ emergent, $SU(3)$ contingent); the multiplicity and $N\le3$ audits; the
internal-3 disposition; the cascade-unification audit; the hostile-review response; and
the v1.0 release. Each of these closures is narrated where it appears later in the
technical parts.

\section{Key discoveries}

\begin{itemize}
\item Synchronization and coherent oscillation emerge spontaneously above a critical
coupling density.
\item Reversibility $\cap$ self-consistency $\,=\,$ quantum mechanics: the Schr\"odinger
equation is the unique maximal-complexity dynamics.
\item $L_Q\to-\nabla^2$ (the elliptic chain) and $\mathrm{BDG}\to\square$ (the Lorentzian
chain) are each controlled continuum limits, but they are \emph{disjoint} in signature.
\item $U(1)$ is a theorem ($\mathrm{Aut}(S^1)=U(1)$, $\pi_1(S^1)=\mathbb Z$), the cleanest
result of the gauge program.
\item The Koide relation is real, predictive, and RG-stable --- and \emph{underivable}
(no-go T-08). The neutrino-Koide analogue is \emph{falsified} (T-11).
\item The RAR $g_\dagger=cH_0/2\pi$ matches $a_0$ with zero free parameters.
\item Complexity maximization selects $d=3{+}1$.
\end{itemize}

\section{Failed routes}

A research program is defined as much by what it abandoned as by what it kept. The
principal abandoned routes were: random matrices as an emergence mechanism; static
topology as a sufficient condition; the apparent resonance-family diversity of TQM-007
(artifacts of a single seed); dynamic graphs as a source of open-ended innovation; the
identification of solitons as \emph{quantum} objects (they are a nonlinear generalization
of carriers, not the quantum regime); open-ended evolution in finite systems (proved to be
delayed saturation); symmetry, attractor, topology, and information-geometry routes to the
Koide relation; topology, attractor, moduli, and persistence routes to the gauge count;
stability, anomaly, representation, defect, and information routes to $N\le3$; tired-light
cosmology; and the TRM Quantum Engine. Each failure was recorded as a no-go or a negative
result, and most were later promoted to the no-go theorems of Chapter~\ref{ch:nogo}.

\chapter{Why the Structure/Content Split}
\label{ch:split}

\section{Historical motivation}

The structure/content split --- \emph{structure is derivable; content is realized} --- was
not an axiom chosen in advance. It was the \emph{conclusion} of a specific audit
(QG-042, the ``Parameter vs.\ Structure'' study) that asked a blunt question: given the
primitives, what can actually be \emph{computed}, and what merely \emph{is}? The answer
divided cleanly along a line nobody had drawn explicitly before.

Everything that could be derived --- oscillation, phase, $U(1)$, charge quantization,
gravity, inertia, the equivalence principle, particle stability, the interpretation of the
Higgs field --- turned out to be \emph{form}, \emph{symmetry}, or \emph{topology}. They
answer \emph{what exists} and \emph{why}. Everything that resisted derivation ---
$\alpha=1/137$, $\alpha_s$, $\theta_W$, the Yukawa couplings, $\lambda$, $\theta_{QCD}$ ---
turned out to be \emph{dimensionless pure numbers}. They answer \emph{how much}. The deep
divide maps exactly onto the two primitives: $Q$ underwrites structure (laws), and Random
Actualization underwrites content (history).

There is a sharp reason the divide is \emph{fundamental} rather than temporary. No
combination of the dimensionful triple $(\ell,\tau,\hbar)$ can produce a dimensionless
number, so a dimensionless coupling \emph{cannot} be assembled from the unit-convention
primitives. The theory therefore does not merely \emph{fail} to derive the couplings; it
explains \emph{why they cannot be derived}: they are historical outcomes of the
actualization draw, not mathematical necessities. Contingency is intrinsic; no multiverse
is needed. Structure is what randomness cannot change; content is what it determines.

\section{What the split is not}

Three confusions must be headed off, because the hostile review raised all three. First,
the split is \emph{not an immunization strategy}: it does not say ``whatever we cannot
derive is contingent''. The boundary is not moved after the fact; it is a fixed, testable
division --- structure is what the primitives compute, content is what they cannot --- and
the theory accepts the consequences, including the underivability of its own couplings.
Second, the split is \emph{not a multiverse}: contingency is intrinsic to a single
universe through Random Actualization, and no ensemble of universes is invoked. Third, the
split is \emph{not anthropic}: the derived pieces (spatial $3$, $U(1)$, the log-normal
form) are theorems or distribution-shape results, not survival-condition selections, and
the selected pieces are flagged as selected rather than hidden behind the word
``derived''.

The point of the split is precisely to make the theory \emph{falsifiable in its
classifications}. A result tagged \realunder\ is a standing invitation to find the
derivation that was missed; a result tagged \drawn\ is a standing prediction that the
value has no deeper origin. Both are testable claims about the boundary of derivability,
not rhetorical immunizations.

\section{Examples}

The split is easiest to grasp through concrete cases, each of which appears in full
later:

\begin{itemize}
\item \textbf{Gauge structure.} $U(1)$ (the group) is \derived; the color count $n=3$
(the content) is \drawn. The group structure is computable from $\mathrm{Aut}(S^1)$; the
count is not fixed by any tested principle (T-09).
\item \textbf{Dimensionality.} Spatial $3$ is \derived (complexity maximization selects
$d=3{+}1$); internal $3$ (generations, color) is \emph{selected} --- a derived lower
bound $N\ge3$ meets an empirical upper bound $N\le3$.
\item \textbf{The abundance law.} The log-normal \emph{form} of the abundance
distribution is \derived (a multiplicative cascade plus the central limit theorem in
log-space); the realized values $\mu,\sigma$ are contingent content.
\item \textbf{The Koide relation.} $Q=2/3$ is real, predictive, and RG-stable --- a
canonical \realunder structure --- but its origin is underivable (T-08).
\item \textbf{Yukawas.} The overlap-operator \emph{form} of the Yukawa operator is
derived; its \emph{spectrum} (the hierarchy $1{:}207{:}3478$) is set by the contingent
architecture shapes.
\end{itemize}

This taxonomy of underivability --- \derived, \realunder, \drawn --- is formalized in
Chapter~\ref{ch:taxonomy}. Its function is honesty: it says precisely which claims are
theorems and which are classifications, and it is the reason the hostile-review history
(Chapter~\ref{ch:audit-trail}) is a history of \emph{framing} corrections rather than of
retracted physics.

\chapter{Hostile Review History}
\label{ch:review-history}

\section{Initial criticisms}

The theory was subjected to two rounds of hostile review, and the record of both is part
of the project. The first review (of the v1.0 paper) raised twelve issues; the response
audit classified them as \emph{1 valid}, \emph{11 partially valid}, and \emph{0 invalid}.
The dominant failure mode was not a fatal theory flaw but a \emph{documentation} gap: the
paper omitted the dynamical system, the formal primitives, the complexity functional, the
emergent-GR matching, and the quantitative predictions that the program actually
contained, and it overstated its framing (``closed'' where it should have said
``dispositioned''; ``no-go'' where it should have said ``conditional no-go''; ``derivation''
where it should have said ``ontological reinterpretation''). Three genuine theory gaps were
acknowledged: the provisional gauge-count no-go (T-09), the contingent content, and the
immunization risk of the structure/content split.

The response audit also identified seven release blockers, all of them documentation fixes
with no new physics: six missing sections (the formal primitive definitions, the
dynamical-system summary, the complexity functional, the emergent-GR derivation summary,
the quantitative predictions, and the scope-and-limitations statement) plus three wording
fixes that were treated as a single framing correction. The important point is the
classification: the review found \emph{one} genuinely valid objection (the ``no route open''
framing), \emph{eleven} partially valid objections that pointed at omitted documentation
rather than at a wrong result, and \emph{zero} fully invalid objections. A review that
lands in that distribution is a review of a paper's \emph{presentation}, not of its
\emph{core} --- and the program treated it accordingly: fix the presentation, keep the
core, and re-verify nothing changed.

A second round (Round~2) raised twelve ``FATAL'' issues, tabulated in
Chapter~\ref{ch:audit-trail}. They ranged from ``primitives undefined as mathematical
objects'' to ``no controlled continuum limit.''

\section{What changed}

The response was a revision that corrected every framing overstatement and added the six
missing sections (formal primitive definitions, dynamical-system summary, complexity
functional, emergent-GR derivation summary, quantitative predictions, scope and
limitations). The wording fixes were: ``closed'' $\to$ ``dispositioned'', ``no-go'' $\to$
``conditional no-go'', and ``derivation'' $\to$ ``ontological reinterpretation'' where
appropriate. No physics changed.

\section{What survived}

The one substantive objection that could not be disposed of by documentation was ``no
controlled continuum limit.'' It was met \emph{fully} on the Schr\"odinger side
($L_Q\to-\nabla^2\to$ Schr\"odinger, plus the curved $L_W$, the Laplace--Beltrami
approximation, and unitarity, all now tested) and \emph{partially} on the Einstein side
(the standard chain works, but the metric and the BDG action are imported). The decisive
residual is therefore the \emph{native metric $\to$ operator coupling} --- the ``G4'' gap:
TQM imports $g_{\mu\nu}$ via Malament's theorem rather than generating it from Q-events.
That residual is what keeps the verdict at \textbf{READY AS A FOUNDATION MONOGRAPH ---
READY AS A RESEARCH-PROGRAM PAPER --- NOT READY AS A DERIVATION PAPER}.

\chapter{The Research Journey}
\label{ch:journey}

The derivation program was organized into research programs, each closed by executable
experiments. Four of them carry the technical weight of this monograph and are narrated
here.

\section{The gauge program}

The gauge program asked why the gauge group is $SU(3)\times SU(2)\times U(1)$. The answer
is tripartite, and each part has a different status.

\textbf{$U(1)$ is derived.} Phase lives on the circle $S^1$; its isometry group \emph{is}
$U(1)$, and the winding $\oint\nabla\theta\cdot dl=2\pi n$ yields integer charge, with the
gauge connection $A_\mu=\partial_\mu\theta$ and the photon as the $n=0$ phase wave. This is
the cleanest result of the entire program: not a choice but the symmetry of phase itself.

\textbf{$SU(2)$ is emergent.} The binary winding pair $\{n=\pm1\}$ is the minimal doublet;
the Bloch sphere gives $SO(3)=SU(2)/Z_2$; the spinor double cover gives $SU(2)$. The
``2'' is near-derived, but the lift from the discrete $Z_2$ to the continuous $SU(2)$ is
not itself derived --- it is the one admitted gap in the $SU(2)$ story.

\textbf{$SU(3)$ is contingent.} The $\mathrm{Aut}(C^n/S_n)\supseteq SU(n)$ argument derives
the group \emph{structure} but not the \emph{count}, and the 8-gluon algebra is borrowed.
Four route families were tested for the count: topology alone (fails for $n>1$, since
$\pi_1(S^1)=\mathbb Z$ is infinite), attractor space (fails, contingent), defect moduli
(the strongest, but it derives structure, not count), and persistence/symmetry (fails to
select). The $1$--$2$--$3$ rank-equals-winding pattern is intriguing but ``possibly
numerology'': no mechanism links the three ranks. The result is no-go T-09 --- no
topological, attractor, moduli, or persistence principle fixes the defect count --- and
the residual ``why 3'' is the theory's one open door.

\subsection{Failed paths and why they failed}

Four route families failed for the defect count. \emph{Topology alone} fails because
$\pi_1(S^1)=\mathbb Z$ is infinite and cannot select $n=3$. \emph{Attractor space} fails
because the landscape content is contingent, so it cannot select a count. \emph{Persistence
and symmetry} fail because every classical group is stable, so stability does not
discriminate. The \emph{rank-equals-winding} ``1--2--3'' pattern is intriguing but has no
linking mechanism --- ``possibly numerology''. The defect-moduli route was \emph{not} a
failure: it is the strongest route, but it derives the structure
$\mathrm{Aut}(C^n/S_n)\supseteq SU(n)$ and stops short of the count, which is exactly the
boundary the no-go (T-09) marks.

\section{The Koide program}

The Koide relation $Q=\sum m/(\sum\sqrt m)^2=2/3$ is real to $\sim10^{-5}$, predictive
(1981 prediction $m_\tau=1776.97$\,MeV, confirmed 1992), and RG-stable. The Koide program
asked where it comes from, and answered by \emph{exhaustion}: it tried every route it
could think of and recorded which failed.

Seven routes were exhausted --- symmetry ($S_3$), topology ($S^1$), attractors,
information geometry, group theory, the $m=3$ closure, and moduli arguments --- with the
result \emph{six no-gos and one blocked}. Four computed origin tests all failed: $S_3$
democratic mixing gives $Q=1/3$ (not $2/3$); there is no RG fixed point at $2/3$; $S^1$
locates but does not fix the angle; and the information-geometry measure
$S/S_{\max}=0.5093$ is not extremal. Five fallacies were explicitly rejected: anthropics
($Q=2/3$ is not anthropic), texture fitting (fits, does not derive), numerology (real but
not derived), hidden parameters (forbidden), and $45^\circ$ restatements (they add zero
information).

The closure is a \emph{no-go} (T-08): the Koide origin is a closed question --- real,
predictive, RG-stable, \emph{underivable} --- the canonical \realunder structure. The
neutrino-Koide analogue was then tested and \emph{falsified} (T-11: $Q_{\max}=0.585$
normal / $0.500$ inverted, both $<2/3$), removing the last remaining distinguisher and
closing the program.

\subsection{Failed paths and why they failed}

Seven routes failed, and the pattern of failure is informative. \emph{Symmetry ($S_3$)}
fails because democratic mixing gives $Q=1/3$ and a hierarchical limit gives $Q=1$; $2/3$
is the non-generic midpoint. \emph{Topology ($S^1$)} fails because it locates the phase
without fixing the $45^\circ$ angle. \emph{Attractors} fail because there is no RG fixed
point at $2/3$. \emph{Information geometry} fails because $S/S_{\max}=0.5093$ is not
extremal. \emph{Group theory} fails to select. The \emph{$m=3$ closure} fails because its
$\gamma\ne2/3$ and its parameters are unmapped to observables. \emph{Moduli} is blocked,
not no-go. The rejected fallacies --- anthropics, texture fitting, numerology, hidden
parameters, $45^\circ$ restatements --- failed because each either fits without deriving
or restates without adding information.

\section{The continuum program}

The continuum program asked whether the discrete dynamical system has a controlled
continuum limit, and it split into two halves with unequal outcomes.

On the Schr\"odinger side the answer is an unambiguous \emph{yes}. The chain Laplacian
$L_Q$ converges to the flat Laplacian $-\nabla^2$ with an exact closed-form spectrum and
second-order convergence, and the weighted Laplacian $L_W$ supplies a curved
generalization and a Laplace--Beltrami approximation. This half is fully tested.

On the Einstein side the answer is more guarded. The BDG causal-set operator converges to
the d'Alembertian $\square$ with $O(h^2)$ control, but $L_Q$ and BDG are \emph{disjoint in
signature} --- one elliptic and positive, the other hyperbolic and indefinite --- so no
native bridge between them has been closed. A dedicated three-test program made the
incompatibility executable: $L_Q$ is positive semi-definite, BDG is indefinite, and the
graph Laplacian is explicitly rejected for Lorentzian signature. This --- two disjoint
chains on the same substrate --- is the single most important structural fact of the
program, and it is the precise reason the Einstein recovery is logical rather than
mathematical.

\subsection{Failed paths and why they failed}

The one path that failed --- and the reason it failed --- is the naive hope that a single
operator could serve both the Schr\"odinger and Einstein sides. It failed because the two
continuum limits have incompatible signatures: $L_Q\to-\nabla^2$ is elliptic and positive
semi-definite, while $\mathrm{BDG}\to\square$ is hyperbolic and indefinite. The executable
test made the failure precise: the graph Laplacian is explicitly rejected for Lorentzian
signature. The program therefore did not produce a missing ``bridge''; it produced a
\emph{no-bridge result} --- two disjoint chains on the same substrate --- which is the
most important structural fact of the program.

\section{The metric program}

The metric program asked how a metric $g_{\mu\nu}$ can arise at all. The answer is a
four-link chain: Q-events $\to$ causal order (native) $\to$ conformal class (proven, via
Malament's imported theorem) $\to$ conformal factor (native, from the counting measure
$f=\rho^{2/d}$) $\to$ the metric. The metric \emph{origin} is therefore closed: the only
imported link is a proven theorem, not a theory gap. The verification is direct --- the
causal order satisfies the metric axioms, the null structure is invariant under conformal
rescaling but not under non-conformal rescaling, and the counting measure round-trips the
conformal factor.

What remains open is the metric \emph{dynamics} --- the native metric $\to$ operator
coupling (G4). Closing the origin tells us \emph{where} the metric comes from; it does not
tell the operator how to \emph{use} it. That residual is why the Einstein recovery remains
logical rather than mathematical, and it is the single genuinely-open technical item in
the monograph.

\subsection{Failed paths and why they failed}

The path that failed is the \emph{native} metric dynamics. TQM does not produce
$g_{\mu\nu}$ from Q-events; the metric-dependent operator $\Delta_g/\Box_g$ is absent, and
the weighted Laplacian $L_W$ supplies only an approximation whose defining metric is
external. The attempt to close the metric \emph{dynamics} natively therefore failed, and
the program's honest endpoint is the split it records: the metric \emph{origin} is closed
(Q-events $\to$ causal order $\to$ conformal class $\to$ conformal factor $\to$ metric,
with Malament's theorem imported as a proven link), while the metric \emph{dynamics} remain
imported. That split is the ``G4'' gap.

\section{What the four programs establish together}

Read together, the four programs trace one arc. The gauge and Koide programs locate the
\emph{boundary} between the derivable and the contingent; the continuum and metric
programs locate the \emph{boundary} between the controlled and the imported. Both
boundaries are honest: the theory claims to derive structure and to control the
Schr\"odinger limit, and it explicitly declines to claim the gauge count, the Koide
origin, or the native Einstein bridge. That division --- derivable structure, controlled
Schr\"odinger limit, classified content, imported metric --- is the shape of the finished
theory.

\chapter{Lessons Learned}
\label{ch:lessons}

\section{Wrong assumptions}

\begin{itemize}
\item \textbf{The static graph was the number-one hidden assumption.} The first
alternative-foundations audit (X001) found eleven hidden assumptions, with the static
graph first on the list. When it was tested (X002--X004), dynamic graphs produced no new
phenomena and no open-ended innovation.
\item \textbf{Nonlinearity breaks the eigenmode picture.} X005 showed that most of the
early ``quantum'' results were linear-algebra artifacts: nonlinearity destroys eigenmodes
and creates solitons instead. Q, fitness, and evolution survived; the naive quantum
reading did not.
\item \textbf{Open-ended evolution in finite systems is an illusion.} X025 appeared to
show open-ended species growth, but X026--X027 proved it was delayed saturation: a
pigeonhole argument bounds the number of species by the number of vertices, so a finite
system always saturates. Level-6 (open-ended) evolution requires an infinite state space.
\item \textbf{Apparent diversity can be a single seed.} The TQM-007 resonance ``families''
collapsed to one universal mode (F4) under a 400-simulation reproducibility study; the
rest were artifacts.
\end{itemize}

\section{Rejected derivations}

The program repeatedly attempted --- and rejected --- candidate derivations rather than
settling for postdiction: symmetry, attractor, topology, and information-geometry routes
to Koide (T-08); topology, attractor, moduli, and persistence routes to the gauge count
(T-09); stability, anomaly, representation, defect, and information routes to $N\le3$
(T-10); five separate attempts to derive $\alpha=1/137$; tired-light cosmology; the TRM
Quantum Engine; and the grand-unified groups $SU(5)/SO(10)$ (proton decay unobserved) and
$E_6/E_8$ (exotic particles unobserved). In each case the rejection was recorded, and
where the failure was systematic it became a no-go theorem.

\section{No-go discoveries}

The no-go theorems (Chapter~\ref{ch:nogo}) are the program's most distinctive scientific
product. T-08 (Koide underivable), T-09 (gauge count unfixed), T-10 ($N\le3$ unfixed),
T-11 (neutrino-Koide falsified), and T-12 (shared cascade untestable) collectively define
the boundary between what the theory can compute and what it must classify. The
meta-lesson is the structure/content split itself: the no-gos are not failures of
ingenuity but the expected signature of a theory in which content is contingent by
construction.

\section{Meta-lessons}

Beyond the specific wrong assumptions and rejected derivations, the program yielded a
small number of standing lessons that shaped everything after them.

\textbf{The derivability boundary is itself a result.} The discovery that structure is
derivable while dimensionless content is not (QG-042) is not a limitation the program
stumbled into; it is a positive finding, and it is what licenses the taxonomy. A theory
that \emph{explains why} its couplings cannot be derived is making a stronger claim than a
theory that merely fails to derive them.

\textbf{Negative results are first-class products.} The no-go theorems (T-08--T-12) are
the program's most distinctive scientific output, not its embarrassments. A route that is
shown to be closed is information, and recording it as a theorem (rather than quietly
dropping it) is what makes the theory auditable.

\textbf{Honesty in framing is a scientific asset.} The hostile-review history is almost
entirely a history of framing corrections, not retracted physics. The lesson is that
``closed'' vs.\ ``dispositioned'', ``no-go'' vs.\ ``conditional no-go'', and ``derivation''
vs.\ ``ontological reinterpretation'' are not cosmetic distinctions; they are the
difference between an overclaim and a defensible statement.

\textbf{Executable verification is what makes honesty checkable.} Every ``PASS'' in this
monograph is the output of a test that a hostile reviewer can run. That discipline is what
turned the hostile-review objections into an audit trail that could be answered point by
point rather than argued about indefinitely.

\chapter{Evolution of the Framework}
\label{ch:evolution}

The framework did not arrive at its final form in one step; it was repeatedly stripped
down and rebuilt. This chapter traces what was assumed at the start, what was removed,
what was retained, and the two methodological commitments that emerged --- the
structure/content split and the no-go method.

\section{Earliest assumptions}

The program began with a rich, and mostly wrong, set of assumptions. Matter was assumed to
be a dynamically stabilized wave structure in a temporal field; synchronization and
resonance were taken to be more fundamental than particles; and the route from oscillators
to cosmology was assumed to be a single continuous chain. Implicitly, the early experiments
also assumed that a static graph was the right substrate, that random matrices might be
sufficient for emergence, and that whatever coherent structure appeared would carry the
quantum reading of the system.

\section{Assumptions removed}

Most of the earliest assumptions were removed by experiment. Random matrices did not
produce emergence (TQM-002); static topology alone was insufficient (TQM-003); the
apparent diversity of resonance families collapsed to a single reproducible mode (TQM-008);
dynamic graphs produced no open-ended innovation (X002--X004); and nonlinearity showed that
the early ``quantum'' results were linear-algebra artifacts (X005). The specific assumptions
of the predecessor TRM framework --- tired-light cosmology and a non-unitary quantum engine
--- were rejected outright in the legacy audit. What remained after this stripping was
minimal: individuation ($Q$), reversibility, self-consistency, and genuine randomness.

\section{Assumptions retained}

Two pairs survived every hostile reduction. \emph{Reversibility} (conserved norm) and
\emph{self-consistency} (fixed points) were proved to be independent and minimally
sufficient: their intersection is unitary quantum mechanics (X011--X015). \emph{$Q$}
(individuation) and \emph{genuine randomness} were shown to be irreducible --- no deeper
invariant exists below them (X035, X039). Everything else in the theory --- time, space,
$3{+}1$, the Schr\"odinger equation, the Born rule, measurement, gravity as causal
structure, particles as defects, and $U(1)$ --- is derived or emergent from these, plus the
single contingent number $M^2$.

\section{Emergence of the structure/content split}

The structure/content split was not imposed; it was \emph{discovered}. The parameter-vs-
structure audit (QG-042) found that everything derivable from the primitives was
\emph{form}, \emph{symmetry}, or \emph{topology}, while everything that resisted derivation
was a \emph{dimensionless pure number}. Since no combination of the dimensionful triple
$(\ell,\tau,\hbar)$ can produce a dimensionless number, the couplings \emph{cannot} be
assembled from the primitives. That turned a temporary limitation into a principle:
structure is derivable, content is realized, and the taxonomy (\derived, \realunder,
\drawn) is the honest accounting of that boundary.

\section{Emergence of the no-go methodology}

The no-go method emerged from the same discipline. When a derivation failed --- Koide, the
gauge count, $N\le3$ --- the program did not discard the failure; it enumerated every route
it could think of, tested each against the no-new-primitives constraint, and recorded the
result as a conditional no-go theorem (T-08--T-12). What began as a way to close individual
questions became a general method: \emph{route exhaustion} plus \emph{hostile reduction},
which turns ``we could not derive it'' into the stronger, checkable statement ``no tested
route derives it.''

\chapterend
{The framework was stripped from a rich oscillator picture to two primitives plus one
number; the structure/content split and the no-go method emerged as its two organizing
commitments.}
{Whether a native metric-to-operator coupling (G4), a unique discriminating prediction, or
a native BDG action can be derived --- the three derivation-paper blockers.}
{Historical/editorial; no executable test applies, but every removal and retention cited
here is recorded in the development ledger (Appendix).}

\chapter{The Method of Executable Audits}
\label{ch:method}

The single most distinctive feature of this program is not a specific result but a
\emph{method}: every substantive claim is made executable, and every ``-audit'' in the
record is a piece of that method. This chapter explains the machinery, because the
confidence statements scattered through the technical parts are unintelligible without it.

\section{What an audit is}

An audit is a structured, reproducible investigation with a single question and a
disposition at the end. It is implemented as an analyzer (a small program) plus one or more
tests, and it writes a report. The question is stated first, the routes are enumerated, and
each route is tested against a fixed set of constraints --- typically the
\emph{no-new-primitives} constraint, and a rejection of anthropics, numerology, hidden
dimensions, and post-selection unless they are explicitly named. The disposition is one of
a small vocabulary: \derived, \realunder, \drawn, \textbf{no-go}, \textbf{blocked},
\textbf{contingent}, or \textbf{selected}. Because the report is written to a file and the
test asserts on the computed numbers, a later reviewer can rerun the audit and get the same
answer.

\section{Route exhaustion and the no-go method}

The no-go theorems are produced by \emph{route exhaustion}: enumerate every derivation
route the program can think of for a target, test each, and record which fail. A no-go is
the statement ``within the tested route space and the no-new-primitives constraint, no
route reaches the target.'' It is therefore \emph{conditional} --- it is not a proof of
logical impossibility (except where a computation, such as T-11, closes the question
directly), and it is only as strong as the enumerated route space. The Koide closure is
the paradigm: seven routes, six no-gos, one blocked, zero open, and the conclusion is a
\emph{closed question}, not a derivation.

\section{Confidence ranks}

The confidence numbers are \emph{uncalibrated ordinal ranks}. They are assigned by the
audit to order the program's own degree of belief in a disposition, and they deliberately
carry no statistical meaning. A theorem with a complete route-exhaustion behind it (such as
$U(1)$, confidence 0.95) ranks above an emergent structure with an admitted gap ($SU(2)$,
0.70), which ranks above a provisional no-go ($SU(3)$ structure, 0.10). The number is a
summary of \emph{how much of the route space has been closed}, not a posterior probability.

\section{Hostile reduction as a method}

The program treats hostile review not as an external event but as a standing internal
method: each result is periodically subjected to a \emph{hostile reduction} that tries to
derive it away, explain it away, or show it is a restatement. The surviving results are
exactly the ones that hostile reduction cannot remove. This is why the hostile-review
history (Chapter~\ref{ch:review-history}) is so short on retracted physics: the framing was
corrected, but the underlying results had already been through hostile reduction
internally before the external referee saw them.

%=============================================================================
\part{Foundations}
%=============================================================================

\chapter{The Primitives}
\label{ch:primitives}

This chapter states the theory. Everything that follows in the technical parts is an
unfolding of these definitions: the primitives are the whole of the input, and the
monograph's burden is to show how much of observable structure follows from them and to
say honestly where it stops.

\section{Statement of the theory}

THE Q-MODEL rests on a single organizational claim, the \textbf{structure/content split}:

\begin{center}
\emph{Structure is derivable; content is realized.}
\end{center}

The \emph{form} of observable structure (topology, symmetry, distribution shape, lower
bounds) is provable from a minimal primitive set. The \emph{content} (specific masses,
couplings, multiplicities, angles) is a draw from a contingent ensemble and is therefore
classified rather than derived.

\section{The four primitives}

\begin{definition}[Primitive set]
The theory admits exactly four primitives, and no others:
\begin{enumerate}
\item \textbf{$Q$} — the principle of individuation (ontology): the irreducible act by
which a discrete event comes to be individuated from nothing.
\item \textbf{Random Actualization} — genuine ontological chance: given $Q$, the realized
event locations are random within the causal structure. (Assumption \textbf{A-03}.)
\item \textbf{$(\ell,\tau,\hbar)$} — the irreducible physical triple: a spacetime scale
$\ell$, a clock $\tau$, and an action $\hbar$ (unit conventions).
\item \textbf{$M^2$} — the single continuous nonlinearity parameter (the nonlinearity
regime of the dynamics).
\end{enumerate}
\end{definition}

No gauge-group, multiplicity, or hidden-dimension primitive is permitted. The primitives
have two layers that meet at $Q$: an \emph{ontology} layer (underwriting the derivation of
structure) and a \emph{dynamical} layer (underwriting the dynamical system of
Chapter~\ref{ch:dynamics}).

\begin{remark}
$Q$ individuates; Random Actualization realizes; $(\ell,\tau,\hbar)$ fixes units; $M^2$
sets the single continuous parameter that the derivation hierarchy pins to $M^2\approx5$
(Chapter~\ref{ch:complexity}).
\end{remark}

\section{The quantum postulates (dynamical layer)}

\begin{axiom}[Q exists]
Topological charge quanta have position $x_i$ and phase $\theta_i$, and interact pairwise
with coupling $J_{ij}=f(|x_i-x_j|)$. (Assumed, irreducible.)
\end{axiom}

\begin{axiom}[Reversible dynamics]
The state norm $\|\psi\|^2$ is conserved; equivalently, evolution is unitary.
\end{axiom}

\begin{axiom}[Born rule]
$P=|\langle\phi|\psi\rangle|^2$, uniquely selected by additivity (Gleason's theorem, 1957).
\end{axiom}

\begin{axiom}[Measurement]
A measurement yields one definite outcome (the collapse axiom).
\end{axiom}

Postulates 1--2 derive the Hilbert space and the Schr\"odinger equation
(Chapter~\ref{ch:dynamics}); postulates 3--4 close the interpretation.

%=============================================================================

\chapterend
{The four primitives ($Q$, Random Actualization, $(\ell,\tau,\hbar)$, $M^2$) and the
structure/content split are stated; no further primitive is permitted.}
{The measure and the action functional for $Q$ remain postulates, not derived objects.}
{No executable test applies to axioms; the formalization audits in the next chapter record
the object/state/operation status.}

\chapter{Formalization of $Q$ and Random Actualization}
\label{ch:formalization}

This chapter records the formal status of the primitives, as established by the
formalization audits (executable; see Part~V). Its purpose is to say, precisely, how much
of each primitive is already a mathematical object and how much remains a postulate ---
the honest boundary that the hostile review pressed on.

\section{Status of $Q$}

\begin{center}
\begin{tabular}{lll}
\toprule
Component & Status & Content \\
\midrule
Object & \textbf{Formalized} & a discrete event (individuated entity) \\
State space & \textbf{Formalized} & positions $x_i$ and phases $\theta_i$ \\
Operations & \textbf{Formalized} & individuation, coupling $J_{ij}$ \\
Configuration space & \textbf{Partial} & graph of events, no continuum limit native \\
Dynamics & \textbf{Partial} & $L_Q$ generator (postulate 1) \\
Symmetries & \textbf{Partial} & $S^1$ phase, $Z_2$, permutation $S_n$ \\
Continuum limit & \textbf{Partial} & verified for $L_Q$ (Part~II) \\
Measure & \textbf{Missing} & no measure space from primitives alone \\
\bottomrule
\end{tabular}
\end{center}

\textbf{Classification:} $Q$ is \emph{partially formalized} — object, state space, and
operations exist; the measure and the action functional remain as postulates.

\section{Status of Random Actualization}

Random Actualization is \textbf{A-03}, an assumption, not a derived object. The
formalization audit found:

\begin{center}
\begin{tabular}{lll}
\toprule
Component & Status & Content \\
\midrule
Random variable & \textbf{Formalized} & the draw of event locations \\
Generator & \textbf{Formalized} & ontological chance \\
$(\Omega,\mathcal F,P)$ & \textbf{Partial} & ensemble unspecified \\
Ensemble measure & \textbf{Partial} & no explicit measure \\
\bottomrule
\end{tabular}
\end{center}

\textbf{Classification:} the random variable and generator are formalized; the probability
space $(\Omega,\mathcal F,P)$ and the ensemble measure are only partially formalized.

\begin{remark}
The hostile reviewer's charge that ``Random Actualization is verbalism'' is therefore
\emph{partially valid}: it is an \emph{admitted} primitive (A-03), not a derived object,
and the paper states this explicitly.
\end{remark}

%=============================================================================

\chapter{The Dynamical System}
\label{ch:dynamics}

This chapter turns the primitives into a dynamics. The single most consequential move is
the identification of the graph Laplacian $L_Q$ as the generator, because its eigenbasis
\emph{is} the Hilbert space of the theory and its continuum limit carries the whole
Schr\"odinger program.

\section{The graph Laplacian $L_Q$}

\begin{definition}[Graph Laplacian]
For the $Q$-interaction graph with adjacency $A$ and degree matrix $D$,
\begin{equation}
L_Q = D - A .
\end{equation}

\eqnotes
{Each row sums to zero, so the uniform state is the zero mode and the eigenvectors are the
standing modes of the graph.}
{It is the discrete diffusion/hopping operator of the theory.}
{Its eigenbasis \emph{is} the Hilbert space, and its continuum limit is the Laplacian that
carries the whole Schr\"odinger program.}
{Undirected and unweighted here; the Lorentzian (BDG) and curved ($L_W$) cases are separate
operators.}
$L_Q$ is real, symmetric, positive semi-definite, and has zero row sums. Its eigenvectors
form an orthonormal basis — the \textbf{Hilbert space} of the theory.
\end{definition}

\section{The tight-binding identity (TQM-142)}

\begin{proposition}[Tight-binding identity]
On a 1D chain,
\begin{equation}
(L_Q)_{ij} = 2\delta_{ij} - \delta_{i,j+1} - \delta_{i,j-1},
\end{equation}
which is \emph{identical} to the tight-binding Hamiltonian with on-site energy
$\varepsilon=2t$: $H = t L_Q$. This is a mathematical \textbf{identity}, not an analogy.
\end{proposition}

\theoremnotes
{On a chain the Laplacian is exactly the second-difference operator, so ``graph diffusion''
and ``tight-binding hopping'' are the same object written twice.}
{The identity is what licenses the transfer of solid-state techniques --- spectra,
Green's functions, band structure --- to the ontology layer without any translation.}
{It holds in one dimension and for the bare (unweighted) Laplacian; the curved and
causal-set generalizations are separate objects (Sections~\ref{ch:laplace-beltrami}
and~\ref{ch:bridge}).}

\section{Schr\"odinger from reversibility (TQM-149--151)}

\begin{theorem}[Schr\"odinger equation from reversibility]
Consider linear evolution $\partial_t\psi = M\psi$ with conserved norm
$\partial_t\|\psi\|^2 = \psi^\dagger(M+M^\dagger)\psi = 0$. Then
\begin{equation}
M^\dagger = -M .
\end{equation}
For real $M$, this is $M^T=-M$. The simplest $2\times2$ antisymmetric matrix $J$ satisfies
$J^2=-I$ (so $J$ acts as the imaginary unit), and
\begin{equation}
M = J\otimes L_Q \quad\Longrightarrow\quad i\partial_t \psi = L_Q\,\psi,
\end{equation}

\eqnotes
{Norm conservation forces the generator to be anti-Hermitian; the block $J$ with $J^2=-I$
supplies the imaginary unit.}
{It is the unitary dynamics of the theory: the Schr\"odinger equation.}
{It means quantum mechanics is \emph{produced} from a conservation requirement, not
postulated.}
{Linear and single-particle; the imaginary structure is imported through $J$.}
with unitary evolution
\begin{equation}
\psi(t) = e^{-iL_Qt}\,\psi(0).
\end{equation}
\end{theorem}

\theoremnotes
{Conservation of the norm forces the generator to be anti-Hermitian; the simplest
anti-Hermitian object on the graph is $J\otimes L_Q$, and $J^2=-I$ is the imaginary unit
in disguise. The ``$i$'' of quantum mechanics is thus not postulated but built from a
$2\times2$ antisymmetric matrix.}
{This is the dynamical core of the theory: the Schr\"odinger equation is obtained from a
conservation requirement plus the graph Laplacian, rather than asserted as an axiom.}
{It is linear and single-particle. The imaginary structure is imported through the
$2\times2$ block $J$, and the argument does not by itself reconstruct the
field-theoretic, interacting content of quantum mechanics.}

\historynotes
{The Schr\"odinger equation is normally \emph{postulated}; the reduction program asked
whether its ``$i$'' and its unitary evolution could instead be derived from a conservation
requirement.}
{If the dynamical core could be derived rather than assumed, the theory would \emph{produce}
quantum mechanics instead of importing it --- a foundational, not cosmetic, gain.}
{The tight-binding/Hilbert-space route used $L_Q$ as a generator, but the specific reduction
of the imaginary unit to a $2\times2$ antisymmetric block $J$ with $J^2=-I$ was the step
that had not been exploited.}
{Accepted as a theorem (TQM-149--151) because it is exact, and because its continuum limit
is then independently verified by \texttt{GraphLaplacianContinuumTests}.}

\section{Observables from the spectrum (TQM-145)}

\begin{center}
\begin{tabular}{ll}
\toprule
Observable & Definition \\
\midrule
effective mass & $m_{\rm eff} = 1/\lambda_1$ \\
energy & $E = \mathrm{tr}(L)$ \\
gap & $\Delta = \lambda_2-\lambda_1$ \\
correlation length & $\xi = 1/\sqrt{\lambda_1}$ \\
diffusion constant & $D = \lambda_1$ \\
channel capacity & $C = \log_2 N$ \\
\bottomrule
\end{tabular}
\end{center}

\begin{remark}
These are \emph{dimensional assignments}, not dynamical predictions; the hostile reviewer
correctly notes this is a scope boundary.
\end{remark}

\section{The ontology $\to$ physics bridge}

The bridge has two legs, both now tested:
\begin{enumerate}
\item $L_Q$ eigenvector basis $=$ Hilbert space (TQM-149); and
\item the causal-set continuum limit yields the metric (QG-001, XC006--012).
\end{enumerate}

%=============================================================================
\part{The Continuum-Limit Program (Verified)}
%=============================================================================

This part records the two continuum chains of the theory and their (partial) connection.
One chain is \emph{elliptic}: the graph Laplacian $L_Q$ converges to the flat Laplacian
$-\nabla^2$, and its weighted generalization supplies a curved-space Schr\"odinger
equation and a Laplace--Beltrami approximation. The other is \emph{Lorentzian}: the
causal-set BDG operator converges to the d'Alembertian $\square$. The two chains are
disjoint in signature, and the part's central fact is that no native bridge between them
has yet been closed. The Schr\"odinger side is controlled and tested; the Einstein side is
partial.

\chapterend
{$L_Q=D-A$ is the generator and its eigenbasis is the Hilbert space; the Schr\"odinger
equation $i\partial_t\psi=L_Q\psi$ follows from norm conservation.}
{The spectrum-derived observables are dimensional assignments, not dynamical predictions.}
{Mathematical (TQM-142 tight-binding identity; TQM-149--151 reversibility theorem); the
continuum-limit tests in the next part verify the downstream chain.}

\chapter{The Two Continuum Chains}
\label{ch:two-chains}

The two chains share a substrate (the Q-events) but use different operators on it. This
chapter exhibits both limits and then states the signature incompatibility that separates
them.

\section{$L_Q \to -\nabla^2$: The Flat Laplacian}
\label{ch:flat-laplacian}

\begin{theorem}[Discrete spectrum of the chain Laplacian]
The graph Laplacian of a 1D chain of $N$ sites has the exact spectrum
\begin{equation}
\lambda_k = \frac{1}{\Delta x^2}\left[2 - 2\cos\frac{\pi k}{N+1}\right],
\qquad k=1,\dots,N,
\end{equation}
with $\Delta x$ the lattice spacing. In the continuum limit $N\to\infty$, $\Delta x\to0$
with the length fixed,
\begin{equation}
\lambda_k \to (\pi k)^2 .
\end{equation}
\end{theorem}

\begin{proof}
The chain Laplacian is the tridiagonal Toeplitz matrix with diagonal $2/\Delta x^2$ and
off-diagonal $-1/\Delta x^2$; its eigenvectors are the discrete sine modes
$\sin(\pi k j/(N+1))$.
\end{proof}

\textbf{Verification.} \texttt{GraphLaplacianContinuumTests} builds chains with
$N=32,64,128,256$, diagonalizes $L_Q$ via MathNet.Numerics, and compares against the closed
form. The relative error \emph{decreases} with $N$ (second-order convergence, error
$\sim4\times$ per doubling), establishing that $L_Q$ converges to the flat Laplacian with
controlled error. \textbf{Result: PASSED.}

\theoremnotes
{Standing waves on a chain are discrete sine modes; taking $N\to\infty$ at fixed length
makes their dispersion relation the usual momentum-squared spectrum of a free particle on
a line.}
{This is the quantitative bridge from the discrete ontology to the continuum: it is an
exact closed form, not a numerical fit, and its convergence is controlled to second order.}
{It is one-dimensional, unweighted, and elliptic only. It says nothing about the
Lorentzian chain or about curvature.}

\historynotes
{The question was whether the discrete ontology has a controlled continuum limit at all ---
the reviewer's ``no controlled continuum limit'' objection.}
{A controlled limit is what separates a toy discrete model from a candidate for continuum
physics; without it the Schr\"odinger program has no bridge to space.}
{Naive finite-difference limits were known, but a closed-form spectrum with a demonstrated
second-order convergence ($\sim4\times$/doubling) had not been established for $L_Q$.}
{Accepted because it is an exact closed form (not a fit) and the error demonstrably
decreases with $N$ in \texttt{GraphLaplacianContinuumTests}.}

\section{BDG $\to \square$: The d'Alembertian}
\label{ch:dalembertian}

The causal-set d'Alembertian is the Benincasa--Dowker--Glaser (BDG) operator: a
directed, alternating sum over causal intervals of the sprinkling.

\begin{theorem}[BDG convergence]
The BDG operator converges to the continuum d'Alembertian $\square$ as the sprinkling
density grows, with leading truncation error $O(h^2)$.
\end{theorem}

\textbf{Verification.} \texttt{BDGOperatorContinuumTests} verifies convergence with a
tolerance $10^{-2}$ (the leading $O(h^2)$ error at $h=1/16$ is $\approx4.37\times10^{-3}$);
the error falls $\sim4\times$ per $h$-halving. \textbf{Result: PASSED.}

\theoremnotes
{A discretized wave operator on a causal set must sum alternatingly over past and future
intervals to cancel the leading terms; the BDG prescription does exactly this, leaving the
continuum d'Alembertian in the dense limit.}
{It establishes the Lorentzian half of the continuum program with controlled $O(h^2)$
error, matching the elliptic half's level of control.}
{The operator is imported from causal-set theory rather than derived from the four
primitives, and the tolerance is leading-order ($10^{-2}$), so only the lowest nontrivial
terms are checked.}

\section{The Quantum-Gravity Bridge}
\label{ch:bridge}

The two continuum chains are \emph{partially connected} — and their signatures are
incompatible.

\begin{proposition}[Two distinct operators]
\begin{enumerate}
\item $L_Q$ is an \emph{undirected, positive} graph Laplacian whose continuum limit is
$-\nabla^2$ (Riemannian, elliptic, positive semi-definite).
\item The BDG operator is a \emph{directed, alternating} causal-set operator whose
continuum limit is $\square$ (Lorentzian, hyperbolic, indefinite).
\end{enumerate}
\end{proposition}

\textbf{Verification.} \texttt{QuantumGravityBridgeTests} (three tests) establish:
\begin{enumerate}
\item $L_Q$ is positive semi-definite (minimum eigenvalue $\ge0$);
\item the BDG operator is indefinite (possesses negative eigenvalues);
\item the two are incompatible in signature — the graph Laplacian cannot serve as a
Lorentzian d'Alembertian (BdgUniquenessAnalyzer O3/O6 explicitly rejects it).
\end{enumerate}
\textbf{Result: PASSED (3/3).}

\begin{remark}
This is the single most important structural fact of the continuum program: TQM possesses
\emph{two} chain types, one Riemannian ($L_Q\to-\nabla^2$) and one Lorentzian
($\mathrm{BDG}\to\square$), which are \emph{disjoint} in signature. The bridge between
them is not yet closed natively.
\end{remark}

\theoremnotes
{A positive semi-definite operator cannot be indefinite, so the elliptic and hyperbolic
continuum limits cannot be the same object; the Q-events underdetermine which operator to
use.}
{This fact is what makes the Einstein recovery \emph{logical} rather than mathematical: the
Schr\"odinger chain and the Einstein chain are separate, and their junction is the open
problem.}
{It states incompatibility, not a construction; it is a no-bridge result, not a bridge.}

%=============================================================================

\chapterend
{$L_Q\to-\nabla^2$ (exact, second-order) and $\mathrm{BDG}\to\square$ ($O(h^2)$) are both
controlled; the two chains are disjoint in signature.}
{No native bridge connects the elliptic and Lorentzian chains.}
{\texttt{GraphLaplacianContinuumTests}, \texttt{BDGOperatorContinuumTests},
\texttt{QuantumGravityBridgeTests} (3) --- all PASS.}

\chapter{The Weighted Laplacian, Laplace--Beltrami, and Curved Schr\"odinger}
\label{ch:weighted}

Having established the flat and Lorentzian chains, this chapter turns to curvature. The
weighted Laplacian $L_W$ is the natural curved generalization of $L_Q$: it supplies a
well-defined curved Schr\"odinger equation and a Laplace--Beltrami approximation, even
though the metric that would make it a genuine $\Delta_g$ remains external.

\section{The weighted Laplacian}

\begin{definition}[Weighted Laplacian]
Given the spatial coupling matrix $K_{ij}$, the weighted graph Laplacian is
\begin{equation}
L_W = D_K - K,
\qquad (D_K)_{ii} = \sum_j K_{ij}.
\end{equation}
This is the standard \emph{unnormalized} graph Laplacian (Belkin--Niyogi), valid for
uniform density.
\end{definition}

\begin{proposition}[Properties]
$L_W$ is symmetric, has zero row sum, is positive semi-definite, and reduces to the
unweighted Laplacian when $K=A$.
\end{proposition}

\textbf{Verification.} \texttt{WeightedLaplacianTests} (four tests) confirm symmetry, zero
row sum, positive semi-definiteness, and reduction to the unweighted case.
\textbf{Result: PASSED (4/4).}

\theoremnotes
{Weighting the edges of the graph lets the geometry be encoded in the couplings; the
unnormalized Laplacian is the operator whose zero-energy state is uniform on the weighted
graph.}
{$L_W$ is the tool that carries the theory from the flat chain to curved space, and its
basic properties are verified rather than assumed.}
{The unnormalized Laplacian is the right operator only at uniform density; at
non-uniform density the normalized variants differ, and none of this supplies the metric
itself.}

\section{The Laplace--Beltrami approximation}
\label{ch:laplace-beltrami}

\begin{theorem}[Laplace--Beltrami from $L_W$]
$L_W$ approximates the Laplace--Beltrami operator $\Delta_g$ in the flat limit and
reproduces standard manifold examples.
\end{theorem}

\textbf{Verification.} \texttt{LaplaceBeltramiTests} (three tests):
\begin{enumerate}
\item \emph{Flat limit}: $L_W$ preserves the flat spectrum;
\item \emph{Variable weights}: non-uniform $K$ changes the spectrum;
\item \emph{Known manifold}: the $S^1$ cycle graph reproduces the circle spectrum.
\end{enumerate}
A subtle degeneracy was found and fixed by audit: the nonzero $S^1$ modes are two-fold
degenerate ($e^{\pm ik\theta}$), so in ascending order the first occurrence of mode $k$ is
at index $2k-1$, not $k$. \textbf{Result: PASSED (3/3).}

\theoremnotes
{The Laplace--Beltrami operator is ``the Laplacian with the metric in it''; approximating
it from a weighted graph means recovering the manifold's spectrum from the couplings.}
{It connects the graph operator to the standard geometric operator of curved space and
exposes (and fixes) the $S^1$ degeneracy subtlety.}
{It is an approximation, not a theorem for arbitrary $g$, and the metric that makes $L_W$
a genuine $\Delta_g$ is still imported rather than generated.}

\section{The curved-space Schr\"odinger equation}
\label{ch:curved-schrodinger}

\begin{theorem}[Curved Schr\"odinger equation]
The operator $L_W$ defines a curved-space Schr\"odinger equation
\begin{equation}
i\,\partial_t \psi = L_W \psi .
\end{equation}
Because $L_W$ is self-adjoint, the evolution is unitary and the norm is conserved:
\begin{equation}
\partial_t \|\psi\|^2 = 0 .
\end{equation}
\end{theorem}

\textbf{Verification.} \texttt{CurvedSchrodingerTests} (three tests): unitarity (norm
conservation), reduction to the flat Schr\"odinger equation, and the existence of a
curved operator. \textbf{Result: PASSED (3/3).}

\begin{remark}
This establishes the \emph{Schr\"odinger} side of the curved-space bridge: the weighted
Laplacian $L_W$ produces a well-defined, unitary curved Schr\"odinger equation. The
\emph{metric} that would make $L_W$ a genuine $\Delta_g$ remains external
(Chapter~\ref{ch:metric-emergence}).
\end{remark}

\theoremnotes
{Self-adjointness is exactly what guarantees unitarity, so the curved equation inherits
norm conservation for free.}
{It closes the Schr\"odinger side of the curved-space program: a well-defined, tested,
unitary curved dynamics exists.}
{The curved equation is formal until a metric is supplied; the metric side is external,
which is the reason the Einstein side remains partial.}

%=============================================================================

\chapterend
{$L_W$ is symmetric/PSD/zero-sum; it gives a Laplace--Beltrami approximation and a unitary
curved Schr\"odinger equation.}
{The metric that would make $L_W$ a genuine $\Delta_g$ remains external.}
{\texttt{WeightedLaplacianTests} (4), \texttt{LaplaceBeltramiTests} (3),
\texttt{CurvedSchrodingerTests} (3) --- all PASS.}

\chapter{The Einstein Tensor Chain}
\label{ch:einstein-tensor}

This chapter exhibits the standard differential-geometry chain from a metric to the
Einstein tensor. The chain is fully functional and tested; the point of including it is to
show \emph{exactly where} it breaks when one tries to feed it TQM's own objects.

\section{Standard differential geometry}

\begin{definition}[Christoffel symbols]
\begin{equation}
\Gamma^\lambda_{\mu\nu} = \tfrac12\, g^{\lambda\sigma}
\left( \partial_\mu g_{\sigma\nu} + \partial_\nu g_{\sigma\mu} - \partial_\sigma g_{\mu\nu} \right).
\end{equation}
\end{definition}

\begin{definition}[Riemann curvature]
\begin{equation}
R^\rho_{\ \sigma\mu\nu} = \partial_\mu \Gamma^\rho_{\nu\sigma}
- \partial_\nu \Gamma^\rho_{\mu\sigma}
+ \Gamma^\rho_{\mu\lambda}\Gamma^\lambda_{\nu\sigma}
- \Gamma^\rho_{\nu\lambda}\Gamma^\lambda_{\mu\sigma}.
\end{equation}
\end{definition}

\begin{definition}[Ricci and Einstein tensors]
\begin{equation}
R_{\sigma\nu} = R^\rho_{\ \sigma\rho\nu}, \qquad
G_{\mu\nu} = R_{\mu\nu} - \tfrac12 R\, g_{\mu\nu}.
\end{equation}
\end{definition}

\section{Verified examples}

\textbf{Unit 2-sphere} $g=\mathrm{diag}(1,\sin^2\theta)$ at $\theta=\pi/4$:
\begin{equation}
\Gamma^\theta_{\phi\phi} = -\sin\theta\cos\theta = -0.5,
\qquad K = 1, \qquad R = 2K = 2, \qquad R_{\mu\nu}=K g_{\mu\nu}.
\end{equation}

\textbf{Unit 3-sphere} $g=\mathrm{diag}(1,\sin^2\chi,\sin^2\chi\sin^2\theta)$ at
$\chi=\theta=\pi/4$: $R_{\mu\nu}=2g_{\mu\nu}$, $R=6$, hence
\begin{equation}
G_{\mu\nu} = R_{\mu\nu} - \tfrac12 R g_{\mu\nu} = -g_{\mu\nu}
= -\mathrm{diag}(1,\tfrac12,\tfrac14).
\end{equation}

\begin{theorem}[2D Einstein tensor vanishes]
In two dimensions $G_{\mu\nu}\equiv0$ identically ($R_{\mu\nu}=\tfrac12 R g_{\mu\nu}$); a
non-trivial $G_{\mu\nu}$ requires dimension $\ge3$.
\end{theorem}

\textbf{Verification.} \texttt{EinsteinTensorTests} (four tests: metric$\to$Christoffel,
Christoffel$\to$Riemann, Riemann$\to$Ricci, Ricci$\to$Einstein) and
\texttt{EinsteinTensorIntegrationTests} (four tests on the integrated builder
\texttt{EinsteinTensorBuilder}). \textbf{Result: PASSED (8/8).}

\theoremnotes
{In two dimensions the Ricci tensor is forced to be proportional to the metric, so the
traceless Einstein tensor vanishes identically; curvature dynamics need at least three
dimensions to have room.}
{It explains, within standard geometry, why a dynamical Einstein equation lives in
$d\ge3$ --- a fact the dimensionality program (Chapter~\ref{ch:complexity}) turns into a
selection argument.}
{It is standard differential geometry, not a TQM derivation; it does not by itself
produce $3{+}1$.}

\begin{remark}[Where the chain breaks in TQM]
The standard chain is fully functional. But TQM's own analyzers
(\texttt{QuantumGravityEmergenceAnalyzer}, \texttt{GrBridgeAnalyzer}) \emph{describe} the
chain as strings (\texttt{GeoStep}) and \emph{never compute} it. The chain breaks at
Step~1 — \emph{metric} $\to$ \emph{Christoffels} — because TQM possesses no native metric
field $g_{\mu\nu}(x)$: it is imported (Chapter~\ref{ch:metric-emergence}).
\end{remark}

%=============================================================================

\chapter{Metric Emergence and Origin Closure}
\label{ch:metric-emergence}

This chapter asks how a metric can arise at all, and records the four-link chain that
closes the metric \emph{origin} --- while leaving the metric \emph{dynamics} imported.

\section{The causal distance structure}

\begin{definition}[Causal volume and distance]
For a causal interval of depth $D$ in $d$ dimensions, the causal volume is
\begin{equation}
N(D) \propto D^d,
\end{equation}
and the causal distance is
\begin{equation}
L = N^{1/d}.
\end{equation}
\end{definition}

\textbf{Verification.} \texttt{MetricGenerationTests} and \texttt{MetricEmergenceTests}
confirm: $L=N^{1/d}$ recovers the depth exactly ($d=2,3,4$; $D=1,\dots,8$); the dimension
is recovered from volume growth; and the resulting distance matrix satisfies all four
metric axioms (non-negativity, identity of indiscernibles, symmetry, triangle inequality).
\textbf{Result: PASSED.}

\section{The conformal factor from the counting measure}

\begin{theorem}[Conformal factor]
The counting measure fixes the volume element $\sqrt{|g|}=\rho$ (density $\rho$). For a
conformally flat ansatz $g=f\,\eta$,
\begin{equation}
\sqrt{|g|} = f^{\,d/2} = \rho \quad\Longrightarrow\quad f = \rho^{\,2/d}.
\end{equation}

\eqnotes
{The event density is the natural volume element, and in a conformally flat ansatz the
volume element fixes the conformal factor.}
{The counting measure determines the one remaining scalar degree of freedom.}
{It closes the conformal \emph{factor} natively, the native half of the metric-origin
chain.}
{Conformally flat only; a general metric is not produced natively.}
\end{theorem}

\textbf{Verification.} \texttt{MetricEmergenceTests} confirms the round-trip
$f^{d/2}=\rho$ for $\rho\in\{0.5,1,2,8\}$, $d\in\{2,3,4\}$. The conformally flat candidate
$g=f\eta$ is constructible: constant $f$ gives $R=0$; $f=1+\tfrac12 x^2$ gives
$R=-0.6145$, matching the Liouville result $K=-\tfrac1{2f}\nabla^2\ln f$.
\textbf{Result: PASSED.}

\theoremnotes
{The density of events is the natural candidate for the volume element, and in a
conformally flat ansatz the volume element alone fixes the conformal factor.}
{This closes the conformal \emph{factor} natively: the counting measure determines the
only remaining scalar degree of freedom.}
{It assumes a conformally flat ansatz; a general metric is not produced natively, only
its conformal class plus a flat representative.}

\section{The conformal class: Malament's theorem}

\begin{theorem}[Malament 1977; Hawking--King--McCarthy 1976]
A bijection preserving the causal order $\prec$ between two past-and-future-distinguishing
spacetimes is a conformal isometry: the causal order determines the metric \emph{up to a
conformal factor}.
\end{theorem}

\textbf{Verification.} \texttt{ConformalStructureTests} and \texttt{MetricOriginTests}
establish: the causal order is a valid partial order whose boundary is the light cone; the
null structure is invariant under $g\to f g$ ($f>0$) but \emph{changed} by a non-conformal
rescale $g=\mathrm{diag}(-1,2)$; hence the light cone picks out \emph{exactly} the
conformal class. \textbf{Result: PASSED.}

\theoremnotes
{Two observers who agree on ``what is before what'' agree on the light cone, hence on the
metric up to an overall local rescale --- the conformal class.}
{This is the imported theorem that turns the causal order into geometry; importing a
proven theorem is standard practice, not a defect.}
{It is imported, not TQM-derived, and it needs past-and-future-distinguishing spacetimes;
it fixes only the conformal class, leaving the factor to be supplied separately.}

\begin{theorem}[Metric origin closure]
The conformal-class ``gap'' is an \emph{imported theorem} (Malament), not a theory gap.
The metric origin is closed:
\begin{equation}
Q\text{-events} \to \text{causal order (native)}
\to \text{conformal class (proven)}
\to \text{conformal factor (native)}
\to g_{\mu\nu}.
\end{equation}
\end{theorem}

\textbf{Classification of the metric origin:}

\begin{center}
\begin{tabular}{lll}
\toprule
Link & Status & Nature \\
\midrule
Q-events & \textbf{NATIVE} & TQM-derived primitive \\
Causal order & \textbf{NATIVE} & precedence from individuation \\
Conformal class & \textbf{IMPORTED} & Malament/HKM — \emph{proven} \\
Conformal factor & \textbf{NATIVE} & counting measure, $f=\rho^{2/d}$ \\
Full $g_{\mu\nu}$ & \textbf{determined} & class $\times$ factor \\
\bottomrule
\end{tabular}
\end{center}

\begin{remark}
The metric origin is \emph{already solved}: importing a proven theorem is standard, not a
defect. The genuinely open item is the \emph{native} metric $\to$ operator coupling
(the ``G4'' gap): TQM does not produce $g_{\mu\nu}$ from Q-events, it imports it.
\end{remark}

\theoremnotes
{The metric is assembled from three native links and one proven import; the ``gap'' is a
proven theorem, so the origin is genuinely closed.}
{It is the theory's answer to ``where does geometry come from'': a chain, not a postulate.}
{Closing the origin does not close the dynamics: the native metric-to-operator coupling
(G4) remains open, which is why Einstein recovery is logical rather than mathematical.}

\historynotes
{The question was whether a metric $g_{\mu\nu}$ can arise from Q-events at all, and where
the theory must stop.}
{Geometry is the substrate of the Einstein side; a native origin would close the largest
remaining gap in the continuum program.}
{Direct construction from Q-events failed (the metric-dependent operator is absent); the
route that worked imports Malament's theorem for the conformal class and supplies the
conformal factor natively.}
{Accepted as closure of the metric \emph{origin} --- the import is a proven theorem, not a
gap --- while the metric \emph{dynamics} (G4) remain honestly imported.}

%=============================================================================
\chapterend
{The metric origin is closed: Q-events $\to$ causal order $\to$ conformal class (Malament,
proven) $\to$ conformal factor ($f=\rho^{2/d}$, native) $\to$ metric.}
{The native metric-to-operator coupling (G4) remains open; metric dynamics are imported.}
{\texttt{MetricGenerationTests} (4), \texttt{MetricEmergenceTests} (4),
\texttt{ConformalStructureTests} (3), \texttt{MetricOriginTests} (3) --- all PASS.}

\part{The Derivation Hierarchy}
%=============================================================================

This part runs the derivation hierarchy from the continuum objects down to concrete
physics: dimensionality, gauge structure, flavor, gravity, and cosmology. Each chapter
records what is \derived, what is \realunder, and what is \drawn, with the confidence
assessments that the hostile-review history refined.

\chapter{Complexity and Spatial Dimensionality}
\label{ch:complexity}

\section{The complexity functional}

The complexity argument is a \textbf{weighted six-component decomposition} of the
conditions for observers (\texttt{ComplexityOptimumAnalyzer.cs}), not a single scalar
variational functional.

\begin{center}
\begin{tabular}{lcc}
\toprule
Component & Weight & $M^2$ window \\
\midrule
Structure formation & 3.0 & 2--8 \\
Particle stability & 2.5 & 3--7 \\
Chemistry potential & 4.0 & 3--5 \\
Information capacity & 2.0 & 3--10 \\
Evolution potential & 1.5 & 2--6 \\
Observer viability & 5.0 & 4--6 \\
\bottomrule
\end{tabular}
\end{center}

\section{Spatial dimensionality $d=3+1$}

\begin{theorem}[Unique dimensionality]
$d=3+1$ is the \emph{unique} value satisfying simultaneously: Bertrand's theorem (closed
orbits only for $1/r^2$), Gauss's law ($F\propto1/r^{d-1}$, Coulomb only in 3D), knot
stability (codimension-2), Huygens (sharp propagation in odd $d$), and 2 GR polarizations.
$d=2+1$ fails gravity/atoms/knots; $d\ge4+1$ fails orbits and knots.
\end{theorem}

\begin{remark}
$M^2\approx5$ is the observer-viability \emph{intersection} of the component windows, not a
hand-inserted number. Confidence T-02 = 0.85; this is a window-intersection argument, not a
variational theorem. The analyzer records that the generation count $G\approx3$ is a
\emph{plateau}, not a peak.
\end{remark}

\theoremnotes
{Each requirement --- closed orbits, Coulomb forces, knots, sharp signals, gravitational
polarizations --- is a filter on dimension, and their intersection is a single value.}
{Dimensionality is derived rather than assumed, closing one of the largest ``why'' questions
with a testable-style intersection argument.}
{It is a window-intersection argument (confidence 0.85), not a variational theorem, and it
derives only the \emph{spacetime} 3; the \emph{internal} 3 (generations, color) is a
separate, selected quantity (Chapter~\ref{ch:nogo}).}

%=============================================================================

\chapterend
{$d=3{+}1$ is the unique value satisfying Bertrand, Gauss, knot stability, Huygens, and the
two GR polarizations simultaneously; $M^2\approx5$ is the observer-viability window.}
{The argument is a window-intersection, not a variational theorem (T-02 $=$ 0.85).}
{Analysis via \texttt{ComplexityOptimumAnalyzer}; not part of the executable test
inventory.}

\chapter{Gauge Structure}
\label{ch:gauge}

This chapter records the gauge program's outcome: $U(1)$ derived, $SU(2)$ emergent, $SU(3)$
contingent. The structure is derivable; the count is not.

\section{$U(1)$ — derived}

\begin{theorem}[$U(1)$ from the circle]
Phase lives on $S^1$. Its isometry group is
\begin{equation}
\mathrm{Aut}(S^1) = U(1), \qquad \pi_1(S^1) = \mathbb{Z},
\end{equation}

\eqnotes
{Phase is an angle, so it lives on a circle; the circle's symmetries are $U(1)$, and going
around it an integer number of times is a conserved charge.}
{$U(1)$ is the symmetry of phase itself.}
{It removes a free gauge assumption from the standard toolkit.}
{Structure only: the Maxwell action is still borrowed, not produced.}
and the winding $\oint\nabla\theta\cdot dl = 2\pi n$ yields integer topological charge.
(Confidence 0.95.)
\end{theorem}

\theoremnotes
{Phase is an angle, so it lives on a circle; the symmetries of a circle are $U(1)$, and
going around the circle an integer number of times is a conserved charge.}
{This is the cleanest result of the gauge program: $U(1)$ is not chosen but is the symmetry
of phase itself, promoted to a theorem (confidence 0.95).}
{It derives the group structure, not the gauge dynamics: the Maxwell action is borrowed, not
produced.}

\historynotes
{The question was whether $U(1)$ electromagnetism is a separate input or follows from the
ontology of phase.}
{$U(1)$ is the oldest gauge symmetry; deriving it from the circle would remove a free
assumption from the standard toolkit.}
{Group-theory approaches \emph{postulate} $U(1)$; the defect/winding route tried to
\emph{derive} it from $\mathrm{Aut}(S^1)$ and $\pi_1(S^1)=\mathbb Z$.}
{Accepted as the cleanest result of the gauge program (confidence 0.95): not a choice but
the symmetry of phase itself --- while admitting the Maxwell action is still borrowed.}

\section{$SU(2)$ and $SU(3)$ — real-underived structure}

\begin{center}
\begin{tabular}{lll}
\toprule
Factor & Status & Confidence \\
\midrule
$U(1)$ & \derived & 0.95 \\
$SU(2)$ & \realunder\ (emergent) & 0.70 \\
$SU(3)$ structure & \realunder\ (emergent) & 0.10 \\
color count $n=3$ & \drawn & — \\
\bottomrule
\end{tabular}
\end{center}

Binary winding $\{n=\pm1\}$ gives the minimal doublet; the Bloch sphere gives
$SO(3)=SU(2)/Z_2$; the spinor double cover gives $SU(2)$. The automorphism of the
$n$-defect moduli space satisfies $\mathrm{Aut}(C^n/S_n)\supseteq SU(n)$, deriving
$SU(3)$ \emph{structure} but not the count.

\begin{remark}
TQM derives the gauge-group \textbf{structure} (which group), \textbf{not} the gauge
\textbf{dynamics} (Maxwell / Yang--Mills actions are borrowed; the 8-gluon algebra is
assumption A-07). The count $n=3$ is unfixed (gauge-count no-go T-09, provisional 0.10).
\end{remark}

%=============================================================================

\chapterend
{$U(1)$ is derived ($\mathrm{Aut}(S^1)=U(1)$, $\pi_1(S^1)=\mathbb Z$); $SU(2)$ is emergent;
$SU(3)$ is contingent (structure, not dynamics).}
{The gauge count $n=3$ is underived (provisional no-go T-09, 0.10).}
{Mathematical (topology); confidence 0.95 / 0.70 / 0.10; no dedicated executable test.}

\chapter{Flavor and the Koide Relation}
\label{ch:flavor}

The flavor chapter is the clearest single illustration of the structure/content split:
the operator \emph{form} is derived, its \emph{spectrum} is contingent, and the Koide
relation sits exactly on the boundary --- real, precise, predictive, and underivable.

\section{The Yukawa operator}

The Yukawa operator is the overlap operator
\begin{equation}
Y_{ij} = \langle \mathrm{arch}_i \,|\,\mathrm{amplitude}\,|\,\mathrm{arch}_j\rangle .
\end{equation}
Its \emph{form} is derived; its \emph{spectrum} (the hierarchy $1{:}207{:}3478$) is
underived — set by the contingent architecture shapes.

\section{The Koide relation}

\begin{theorem}[Koide relation]
The charged-lepton pole masses satisfy
\begin{equation}
Q = \frac{\sum m}{(\sum\sqrt{m})^2} = 0.6666605,
\qquad \theta = 44.9997^\circ,
\end{equation}

\eqnotes
{In $\sqrt m$ (amplitude) space the three charged leptons sit at $45^\circ$ to the
democratic direction.}
{$Q=2/3$ is $1/(3\cos^2 45^\circ)$; the relation is about amplitudes, not masses.}
{It is the sharpest numerical clue in the flavor sector: real, predictive, RG-stable.}
{Its origin is underivable (no-go T-08): real does not mean derived.}
equivalently $Q=2/3$ at $\theta=\arccos(1/\sqrt2)=45^\circ$. The relation is real
($\sim10^{-5}$ precision; Bayes factor vs.\ coincidence $\approx3.2\times10^4$), predictive
(1981 prediction $m_\tau=1776.97$ MeV, confirmed 1992), and RG-stable.
\end{theorem}

\theoremnotes
{In $\sqrt m$ (i.e., amplitude) space the three charged leptons sit at a $45^\circ$ angle to
the democratic direction, which is exactly the condition $Q=2/3$; the relation is about
amplitudes, not masses.}
{It is a real, predictive, RG-stable regularity --- the strongest numerical clue in the
flavor sector --- and the canonical example of a \realunder structure.}
{Its origin is underivable (no-go T-08 below): real does not mean derived, and the relation
does not by itself fix the third-generation count.}

\historynotes
{The Koide relation has stood since 1981 as an unexplained, high-precision regularity; the
program asked whether its origin is derivable from the primitives.}
{It is the single sharpest numerical clue in the flavor sector; its origin would either
confirm a deep structure or prove that structure stops at form.}
{Seven routes were tried --- symmetry, topology, attractors, information geometry, group
theory, $m=3$ closure, moduli --- and all failed or blocked.}
{Accepted as a \emph{closed question} (no-go T-08): real, predictive, RG-stable, and
underivable --- the canonical \realunder structure, with the neutrino-Koide analogue
falsified (T-11).}

\begin{nogo}[Koide origin — T-08]
No symmetry ($S_3$), attractor, topology ($S^1$), or information-geometry selects $Q=2/3$.
(Confidence 0.70.)
\end{nogo}

The closure audit (Phase 159) exhausted seven routes:

\begin{center}
\begin{tabular}{ll}
\toprule
Route & Status \\
\midrule
Symmetry ($S_3$) & No-Go \\
Topology ($S^1$) & No-Go \\
Attractors & No-Go \\
Information Geometry & No-Go \\
Group Theory & No-Go \\
$m=3$ Closure & No-Go \\
Moduli Arguments & Blocked \\
\bottomrule
\end{tabular}
\end{center}

The Koide origin is a \emph{closed question} — real, predictive, RG-stable,
\emph{underivable} — the canonical \realunder\ structure.

\begin{nogo}[Neutrino-Koide — T-11, falsification]
Requiring $Q=2/3$ for neutrinos gives $Q_{\max}=0.585$ (normal) / $0.500$ (inverted), both
$<2/3$; the measured $\Delta m^2$ exclude all-lepton Koide. (Confidence 0.90.)
\end{nogo}

%=============================================================================

\chapterend
{The Koide relation is real, predictive, and RG-stable; its origin is underivable (T-08);
the neutrino-Koide analogue is falsified (T-11).}
{None: the Koide origin is a closed question, and the falsification closed the last
distinguisher.}
{Computational audits ($Q$, $Q_{\max}$); recorded in the Koide closure, not in the
executable test inventory.}

\chapter{Gravity and Emergent GR}
\label{ch:gravity}

Gravity in TQM is a phase-gradient phenomenon. This chapter records the honest status:
the Einstein recovery is \emph{logical}, not a new dynamical derivation.

\section{Newton's constant}

\begin{proposition}[Dimensional result]
$G = \ell^2 c^3/\hbar$ is a \emph{unit-consistency} result, not a derivation of gravity's
dynamics.
\end{proposition}

\eqnotes
{The triple $(\ell,\tau,\hbar)$ can assemble exactly one combination with the dimensions of
$G$.}
{$G$ is fixed by units once the scale $\ell$ is fixed.}
{It removes Newton's constant from the list of independent inputs.}
{Dimensional analysis only; it says nothing about the field equations.}

\theoremnotes
{The triple $(\ell,\tau,\hbar)$ can assemble exactly one combination with the dimensions of
$G$, so Newton's constant is fixed by units once the scale is fixed.}
{It removes $G$ from the list of independent constants --- gravity's strength is a
bookkeeping consequence of the scale triple.}
{It is dimensional analysis, not dynamics; it says nothing about the field equations.}

\section{The phase-gravity chain}

oscillation $\to$ phase $\to$ causal density $\to$ metric $\to$ curvature $\to$ Einstein.

\begin{theorem}[Leading-order recovery]
\begin{equation}
G_{\mu\nu} = 8\pi G_{\rm eff}\, T_{\mu\nu} + O(\ell_P^2 R^2),
\end{equation}
with Planck-scale corrections $\sim10^{-40}$. Newtonian $1/r^2$, lensing, redshift,
perihelion precession, and two gravitational-wave polarizations at speed $c$ are reproduced.
\end{theorem}

Two modifications beyond GR: (1) time-varying dark energy
\begin{equation}
\Lambda(t) = \frac{\alpha}{\sqrt{V(t)}},
\end{equation}
where $V(t)$ is the 4-volume of the past light cone; and (2) singularity resolution at
$r\sim\ell_P$ (maximum curvature $1/\ell_P^2$).

\begin{remark}[Honesty note]
The chain is \textbf{logical, not a new dynamical derivation}: ``same equations, same
predictions; no falsifiable difference'' in the weak field. The value is \emph{ontological}
(identifying the gravitational potential with the phase field); the physical content beyond
GR is confined to $\Lambda(t)$ and singularity resolution.
\end{remark}

\theoremnotes
{If the phase field is the gravitational potential, then the standard curvature chain is
recovered to leading order, with only Planck-scale corrections --- gravity is re-labeled,
not replaced.}
{It shows the ontological fit is consistent with GR in the weak field while isolating the
two places where TQM genuinely departs ($\Lambda(t)$ and singularity resolution).}
{It is logical, not mathematical: the metric and BDG action are imported, so this is a
reinterpretation, not a derivation of Einstein's equations.}

\section{The radial acceleration relation}

\begin{theorem}[Zero-parameter RAR]
\begin{equation}
g_\dagger = \frac{c H_0}{2\pi} \approx 1.05\times10^{-10}\ \mathrm{m/s^2},
\end{equation}

\eqnotes
{$cH_0$ is the natural acceleration scale of a causal set; dividing by $2\pi$ reproduces the
Milgrom scale $a_0$.}
{It is a zero-parameter prediction that lands on the measured value.}
{It is the theory's sharpest dark-sector result.}
{The $2\pi$ is an admitted numerical accident; this is a match, not yet a full derivation.}
matching the measured $a_0$. (The factor $2\pi$ is a numerical accident, Phase 144.)
\end{theorem}

\theoremnotes
{The combination $cH_0$ is the natural acceleration scale of a causal set; dividing by $2\pi$
reproduces the observed Milgrom scale $a_0$.}
{It is a zero-parameter prediction that lands on the measured value, a genuine sharp result
in the dark-sector.}
{The $2\pi$ factor is a numerical accident, and the relation is a match, not yet a full
derivation of the dynamics behind $a_0$.}

%=============================================================================

\chapterend
{$G=\ell^2c^3/\hbar$ is dimensional; the leading-order Einstein recovery is logical; the
RAR $g_\dagger=cH_0/2\pi$ matches $a_0$ with zero free parameters.}
{Native metric dynamics (G4); a unique, sharp discriminating prediction.}
{\texttt{EinsteinTensorTests} (4) and \texttt{EinsteinTensorIntegrationTests} (4) verify
the standard chain; \texttt{EinsteinRecoveryTests} (3) = PARTIAL.}

\chapter{Cosmology}
\label{ch:cosmology}

Cosmology is the least closed chapter of the monograph: it is an accepted partial
computational layer, not a completed derivation. The status is recorded without
overclaim.

\begin{itemize}
\item \textbf{Expansion} is an interpretation of FLRW (QG-081); every reinterpretation
collapses to FLRW (QG-080--089).
\item \textbf{Causal-set $\Lambda$}: $\Lambda\sim1/\sqrt{N}$ is a postdiction (Phase 140).
\item \textbf{Pantheon+SH0ES}: TQM and $\Lambda$CDM are indistinguishable at the current
signal $w(z)=-1+0.015(1+z)^{3/2}$ (DATA-001); a detectability audit (DATA-002) quantified
the power required for separation.
\item \textbf{Clusters}: Coma dynamical mass and ACCEPT X-ray gas fraction are reproduced.
\item \textbf{CMB}: an accepted partial computational layer; full $C_\ell$ solver deferred
(computational, not a theory gap).
\end{itemize}

\begin{remark}
The cosmology status is deliberately partial: the background and the compression peaks of
the CMB are present, but the rarefaction peak and the acoustic phase shift require a
CAMB-class Boltzmann solver --- standard $\Lambda$CDM physics, not new TQM physics.
\end{remark}

%=============================================================================
\chapterend
{Expansion is an FLRW reinterpretation; causal-set $\Lambda\sim1/\sqrt N$ is a postdiction;
the $w(z)$ signal is quantified.}
{The full CMB $C_\ell$ solver is deferred (computational, not a theory gap).}
{Partial computational layer (DATA-001/002, CMB closure audits); not in the executable test
inventory.}

\part{Classification, No-Gos, and Predictions}
%=============================================================================

\chapter{The Three-Category Taxonomy}
\label{ch:taxonomy}

The taxonomy is the theory's honest-accounting device. It classifies \emph{underivability},
so that no claim is allowed to pass as a derivation when it is a classification.

\begin{definition}[Taxonomy]
\begin{center}
\begin{tabular}{ll}
\toprule
Category & Meaning \\
\midrule
\derived & computable from the primitives by theorem \\
\realunder & real, precise, predictive structure whose origin is not computable \\
\drawn & a coincidental draw of the abundance law \\
\bottomrule
\end{tabular}
\end{center}
The taxonomy classifies \emph{underivability}; \realunder\ and \drawn\ are \emph{not}
claimed as derivations.
\end{definition}

\section{The fourteen classified results}

\begin{center}
\begin{tabular}{llc}
\toprule
Object & Category & Confidence \\
\midrule
$U(1)$ & \derived & 0.95 \\
spatial 3 & \derived & 0.85 \\
$N\ge3$ & \derived & 0.90 \\
log-normal law & \derived & theorem \\
$SU(2)$ & \realunder\ (emergent) & 0.70 \\
$SU(3)$ structure & \realunder\ (emergent) & 0.10 \\
Koide $Q=2/3$ (reality) & \realunder\ (structured) & 0.90 \\
Koide $45^\circ$ (origin) & \realunder\ (structured) & 0.70 \\
Yukawas / couplings / $\Omega_{DM}$ & \drawn & — \\
$N\le3$ & \drawn & 0.70 \\
color count 3 & \drawn & — \\
internal $N=3$ & \derived$\cap$\drawn & 0.70 \\
neutrino-Koide & FALSIFIED & 0.90 \\
\bottomrule
\end{tabular}
\end{center}

Category conflicts: \textbf{0}. Composite objects: \textbf{2}. Overall consistency:
\textbf{0.95}.

%=============================================================================

\chapter{The No-Go Theorems}
\label{ch:nogo}

The no-go theorems are \emph{conditional} — relative to the no-new-primitives constraint and
the tested route space — not absolute logical impossibilities (except the computation
T-11).

\begin{center}
\begin{tabular}{p{3.4cm}p{9.6cm}c}
\toprule
Theorem & Statement & Confidence \\
\midrule
\textbf{T-08} Koide no-go & no symmetry/attractor/topology/information-geometry selects $Q=2/3$ & 0.70 \\
\textbf{T-09} Gauge-count no-go (provisional) & no principle fixes $n$; graph-spectrum/lattice-mode untested & 0.10 \\
\textbf{T-10} $N\le3$ no-go & no stability/anomaly/representation bound gives $N\le3$ & 0.70 \\
\textbf{T-11} Neutrino-Koide falsification & $Q_{\max}=0.585/0.500 < 2/3$ (computation) & 0.90 \\
\textbf{T-12} Shared-cascade no-go & one vs three cascades untestable from one universe & 0.55 \\
\bottomrule
\end{tabular}
\end{center}

\begin{remark}[Internal-3 disposition]
Gauge count $n=3$ and multiplicity $N\le3$ are one node with two formally unlinked faces
($N$ vs $n$; assumption A-10). It is dispositioned \emph{unresolved-contingent}: the
multiplicity face closed (T-10, 0.70); the gauge-count face \emph{provisionally} closed
(T-09, 0.10). This is the theory's one open door.
\end{remark}

\theoremnotes
{Each no-go is a statement about what the tested route space cannot reach, not about what
is logically impossible; the ``3'' that pervades physics reduces to one node with two
faces.}
{The no-gos are the program's most distinctive product: they locate the exact boundary
between the derivable and the contingent, which is the content of the structure/content
split.}
{They are conditional (relative to no-new-primitives and the tested routes), the gauge-count
no-go T-09 is only provisional (0.10), and the shared-cascade no-go T-12 is weak (0.55).}

%=============================================================================

\chapterend
{T-08--T-12 are conditional no-gos; the internal-3 node is the theory's one open door.}
{The gauge-count face (T-09, 0.10) is only provisionally closed.}
{T-11 is a computation (falsification); T-08 is route exhaustion; the rest are recorded
audits.}

\chapter{Quantitative Predictions}
\label{ch:predictions}

The theory's claim to be scientific rests on the fact that it makes predictions that can
fail, and has already failed one. The table records the live predictions and the one
falsified prediction; the falsification is kept in the table rather than buried, because
it is the strongest evidence that the program is not immunized.

\begin{center}
\begin{tabular}{lll}
\toprule
Prediction & Type & Status \\
\midrule
RAR $g_\dagger=cH_0/(2\pi)\approx1.05\times10^{-10}$ & zero-parameter & matches $a_0$ \\
$w(z)=-1+0.015(1+z)^{3/2}$ & cosmological & Euclid $>3\sigma$ by $\sim$2030 \\
$\Lambda(t)=\alpha/\sqrt{V(t)}$ & dark energy & Euclid / Roman \\
log-normal abundance law & distribution form & testable \\
$N\ge3$ (CP lower bound) & a priori & confirmed \\
neutrino-Koide $Q=2/3$ & falsifiable & \textbf{falsified} \\
\bottomrule
\end{tabular}
\end{center}

The theory is \emph{falsifiable}: it made a concrete prediction (neutrino-Koide) that was
falsified, and it carries live quantitative predictions that distinguish it from SM
$+\Lambda$CDM in the \emph{form} sector, even though the \emph{content} is DRAWN.

%=============================================================================

\chapter{Scope and Limitations}
\label{ch:scope}

A theory that states its limitations is more useful than one that hides them. The
following list is the theory's own accounting of where it stops; it is reproduced from the
hostile-review response and is intended to be read as part of the theory, not as an
apology for it.

\begin{enumerate}
\item The gauge-count face of the internal-3 node is only provisionally closed (T-09, 0.10).
\item Encyclopedia completeness $\approx81\%$; the CMB chapter is a documented partial
computational layer.
\item Content is contingent by design; the theory predicts the distribution \emph{shape},
not the \emph{values}.
\item The shared-cascade question is logically (not empirically) closed.
\item The unified action is a roadmap, not a result.
\item Confidence numbers are uncalibrated ordinal ranks (Phase-149/156 estimates).
\item The phase-gravity chain is ontological in the weak field.
\end{enumerate}

%=============================================================================
\part{Verification}
%=============================================================================

\chapter{The Executable Verification Program}
\label{ch:verification}

The claim that this is a \emph{tested} theory rests on the executable program summarized
here and inventoried in full in Appendix~\ref{app:tests}. Every ``PASS'' below is the
output of an xUnit test, not a hand computation.

\section{Summary of verified results}

\begin{center}
\begin{tabular}{lll}
\toprule
Chain link & Test file & Result \\
\midrule
$L_Q\to-\nabla^2$ & GraphLaplacianContinuumTests & exact spectrum, 4$\times$/doubling \\
BDG$\to\square$ & BDGOperatorContinuumTests & $O(h^2)$, 4$\times$/halving \\
$L_Q$ vs BDG signature & QuantumGravityBridgeTests (3) & PSD vs indefinite \\
no native $\Delta_g$ & CurvedSpaceBridgeTests (3) & bridge absent \\
weighted Laplacian & WeightedLaplacianTests (4) & symmetric/PSD/zero-sum \\
$L_W\to\Delta_g$ & LaplaceBeltramiTests (3) & S¹ example \\
curved Schr\"odinger & CurvedSchrodingerTests (3) & unitary \\
Einstein recovery & EinsteinRecoveryTests (3) & PARTIAL \\
Einstein chain & EinsteinTensorTests (4) & standard 2D/3D \\
Einstein integration & EinsteinTensorIntegrationTests (4) & builder works \\
metric generation & MetricGenerationTests (4) & distance PRESENT \\
metric emergence & MetricEmergenceTests (4) & factor native \\
conformal structure & ConformalStructureTests (3) & class imported \\
metric origin & MetricOriginTests (3) & closure \\
\bottomrule
\end{tabular}
\end{center}

\section{The decisive findings}

\begin{enumerate}
\item The \emph{Schr\"odinger} side of the continuum limit is \textbf{controlled and
tested}: $L_Q\to-\nabla^2\to$ Schr\"odinger (and the curved $L_W$, and the
Laplace--Beltrami approximation) are verified.
\item The \emph{Einstein} side is \textbf{partially} verified: the standard $g\to
G_{\mu\nu}$ chain works, but TQM has no native metric (Malament import), and the BDG action
is imported.
\item The conformal-class ``gap'' is an \emph{imported theorem} (Malament), not a theory
gap.
\end{enumerate}

%=============================================================================

\chapterend
{The Schr\"odinger side of the continuum limit is controlled and tested; the Einstein side
is partial; the conformal class is an imported proven theorem.}
{The native metric-to-operator coupling (G4).}
{15 test files / 47 tests, all PASS (\texttt{EinsteinRecoveryTests} noted PARTIAL).}

\chapter{The Hostile-Review Audit Trail}
\label{ch:audit-trail}

This chapter is the tabular record of the hostile-review history narrated in
Chapter~\ref{ch:review-history}. It lists the Round-2 ``FATAL'' issues and their final
classifications.

A hostile referee (Round~2) raised twelve FATAL issues. Each was re-evaluated against the
executable tests. The classifications are:

\begin{center}
\begin{tabular}{p{6.5cm}l}
\toprule
Round-2 FATAL issue & Classification \\
\midrule
Primitives undefined as mathematical objects & PARTIALLY RESOLVED \\
Schr\"odinger derivation circular/incomplete & RESOLVED \\
Gauge derivations = elementary topology & PARTIALLY RESOLVED \\
Complexity argument not a derivation & PARTIALLY RESOLVED \\
Structure/content split = immunization & PARTIALLY RESOLVED \\
Composite objects admitted & PARTIALLY RESOLVED \\
Internal-3 unresolved while ``complete'' & RESOLVED \\
T-09 at 0.10 cannot close & RESOLVED \\
``Structurally complete'' overstatement & RESOLVED \\
Gravity$\to$Einstein is ontological & RESOLVED \\
$\mathrm{Aut}(S^1)=U(1)$ as EM derivation & RESOLVED \\
No controlled continuum limit & PARTIALLY RESOLVED \\
\bottomrule
\end{tabular}
\end{center}

\textbf{Tally: RESOLVED = 6 · PARTIALLY RESOLVED = 6 · OPEN = 0.}

\section{What changed}

The framing overstatements (``closed'' vs.\ ``dispositioned'', ``no-go'' vs.\
``conditional no-go'', ``derivation'' vs.\ ``ontological reinterpretation'') were
corrected in the revised paper. The \emph{substantive} objection — ``no controlled
continuum limit'' — was met on the Schr\"odinger side (now tested: $L_Q\to-\nabla^2\to$
Schr\"odinger, and curved $L_W$, Laplace--Beltrami, unitary), and \emph{partially} on the
Einstein side (the standard chain works, but the metric and BDG action are imported).

\section{The decisive residual}

The single genuinely-open item is the \textbf{native metric $\to$ operator coupling} (the
``G4'' gap): TQM imports $g_{\mu\nu}$ (Malament) rather than generating it from Q-events.
The metric \emph{origin} is closed (Chapter~\ref{ch:metric-emergence}); the metric
\emph{dynamics} remain imported.

\textbf{Verdict:} READY AS A FOUNDATION MONOGRAPH --- READY AS A RESEARCH-PROGRAM PAPER
--- NOT READY AS A DERIVATION PAPER (see Chapter~\ref{ch:conclusion}).

%=============================================================================

\chapter{What Survived Hostile Review}
\label{ch:survived}

This chapter consolidates the two hostile-review rounds into a single survival table. For
each major claim, it records the original status, the reviewer's criticism, and the final
status after re-evaluation against the completed tests. The pattern is uniform: nothing
was retracted, but many overstatements were downgraded to precise, defensible forms.

\begin{longtable}{p{3.2cm}p{2.4cm}p{3.0cm}p{5.8cm}}
\toprule
Claim & Original status & Review criticism & Final status \\
\midrule
\endhead
Primitives & ``fully formalized'' & undefined as mathematical objects & partially formalized (measure/action admitted missing) \\
Schr\"odinger equation & derived & circular / incomplete & derived and tested ($L_Q\to-\nabla^2\to$ Schr\"odinger, unitary) \\
Gauge group & derived & elementary topology & $U(1)$ derived; $SU(2)$ emergent; $SU(3)$ contingent (structure, not dynamics) \\
Dimensionality $3{+}1$ & derivation & not a derivation & window-intersection (T-02 $=$ 0.85) \\
Structure/content split & organizing principle & immunization & retained; falsifiability demonstrated (neutrino-Koide) \\
``Structurally complete'' & complete & overstatement & ``dispositioned''; internal-3 flagged open \\
No-go theorems & no-go & too strong & conditional no-go (relative to tested route space) \\
Gravity $\to$ Einstein & derivation & ontological, not derivation & ontological reinterpretation in the weak field \\
$G=\ell^2c^3/\hbar$ & derived & dimensional analysis & dimensional / unit-consistency result \\
Continuum limit & controlled & no controlled limit & Schr\"odinger side controlled and tested; Einstein side partial \\
Metric origin & closed & a gap & origin closed (Malament import, proven); dynamics imported (G4) \\
Unique prediction & testable & no discriminating prediction & live predictions; none yet unique vs.\ SM$+\Lambda$CDM \\
Koide relation & real structure & numerology & real, predictive, RG-stable; origin underivable (T-08) \\
Gauge count $n=3$ & derived & unfixed & provisional no-go T-09 (0.10) \\
\bottomrule
\end{longtable}

\textbf{Reading of the table.} Every ``derived'' claim that survived hostile review became
either a tested theorem ($U(1)$, the Schr\"odinger limit) or an explicitly scoped statement
(window-intersection, ontological reinterpretation, unit-consistency). Every ``gap'' alleged
by the reviewers resolved into one of two honest categories: an imported proven theorem
(Malament) or an admitted theory limit (G4, the gauge count, the discriminating
prediction). Nothing in this table was retracted, and nothing was silently dropped.

\chapterend
{No claim was retracted; overstatements were downgraded to precise forms; the single
genuine theory limit (G4) and the three derivation-paper blockers are now explicit.}
{The same three scientific blockers (native metric coupling, unique prediction, native BDG
action), which affect only the derivation-paper claim.}
{Tabulated from the completed tests and the two hostile-review rounds; no executable test
applies to the table itself.}

\chapter{Conclusion}
\label{ch:conclusion}

THE Q-MODEL is \emph{structurally complete} in the precise sense that every in-scope
question is \textbf{dispositioned} — derived, underivable, underdetermined, contingent, or
accepted as a computational layer.

The single residual of genuine scientific tension is the \textbf{internal-3 node} — why the
internal multiplicity/count saturates at 3 — dispositioned unresolved-contingent with a
provisional gauge-count no-go (T-09, 0.10). This is the theory's one open door.

\textbf{Structure is derivable. Content is realized.}

\textbf{Publication status:} READY AS A FOUNDATION MONOGRAPH --- READY AS A
RESEARCH-PROGRAM PAPER --- NOT READY AS A DERIVATION PAPER. The Schr\"odinger dynamical
core is executable and tested; the Einstein recovery remains logical, not mathematical
(imported metric + imported BDG action, $G$ dimensional, no unique discriminating
prediction). The foundation-monograph and research-program claims are unaffected by the
remaining derivation-paper blockers, because a foundation monograph may legitimately use
imported proven theorems.

%=============================================================================
\appendix
%=============================================================================

\chapter{Complete Test Inventory}
\label{app:tests}

\begin{longtable}{lllc}
\toprule
Test file & Tests & Verified quantity & Status \\
\midrule
\endhead
GraphLaplacianContinuumTests & 1 & $\lambda_k=(1/\Delta x^2)[2-2\cos(\pi k/(N{+}1))]\to(\pi k)^2$ & PASS \\
BDGOperatorContinuumTests & 1 & BDG$\to\square$, $O(h^2)$ & PASS \\
QuantumGravityBridgeTests & 3 & $L_Q$ PSD; BDG indefinite; incompatible & PASS \\
CurvedSpaceBridgeTests & 3 & no $\Delta_g$; $\Delta_g\to\nabla^2$; absent & PASS \\
MetricOperatorTests & 4 & candidate constructions & PASS \\
WeightedLaplacianTests & 4 & symmetric; zero-sum; PSD; unweighted limit & PASS \\
LaplaceBeltramiTests & 3 & flat limit; weights; S¹ ($2k{-}1$ degeneracy) & PASS \\
CurvedSchrodingerTests & 3 & unitary; flat limit; curved operator & PASS \\
EinsteinRecoveryTests & 3 & PARTIAL (no full $G_{\mu\nu}$) & PASS \\
EinsteinTensorTests & 4 & flat/sph $\Gamma$, $K$, $R$, $G$ & PASS \\
EinsteinTensorIntegrationTests & 4 & builder; 3-sphere $G=-g$ & PASS \\
MetricGenerationTests & 4 & distance PRESENT; candidate PARTIAL & PASS \\
MetricEmergenceTests & 4 & factor $f=\rho^{2/d}$; $R=-0.6145$ & PASS \\
ConformalStructureTests & 3 & axioms; null invariant; Malament & PASS \\
MetricOriginTests & 3 & class unique; factor free; imported & PASS \\
\bottomrule
\end{longtable}

\chapter{Confidence and Classification Tables}
\label{app:confidence}

The confidence numbers are uncalibrated ordinal ranks (Phase-149/156 estimates), not
posterior probabilities. They rank, they do not measure.

\section{Confidence table}

\begin{center}
\begin{tabular}{llc}
\toprule
Result & Confidence & Note \\
\midrule
program consistency & 0.95 & zero category conflicts \\
$U(1)$ & 0.95 & theorem \\
$N\ge3$ & 0.90 & CP lower bound \\
neutrino-Koide & 0.90 & falsification (computation) \\
spatial 3 & 0.85 & window intersection \\
$SU(2)$ & 0.70 & emergent \\
Koide $45^\circ$ & 0.70 & structured \\
$N\le3$ & 0.70 & no-go \\
$SU(3)$ structure & 0.10 & provisional no-go \\
\bottomrule
\end{tabular}
\end{center}

\section{Final classification}

\begin{center}
\begin{tabular}{ll}
\toprule
Question & Classification \\
\midrule
Is the metric origin closed? & YES (imported Malament, proven) \\
Is the Schr\"odinger continuum limit controlled? & YES (tested) \\
Is the Einstein recovery a derivation? & NO (logical, imported) \\
Is TQM journal-ready? & NO \\
Is TQM white-paper-ready? & YES \\
\bottomrule
\end{tabular}
\end{center}

%=============================================================================

\chapter{Worked Derivation: The Einstein Chain on the 2-Sphere}
\label{app:worked}

This appendix exhibits the standard differential-geometry chain explicitly, so the reader
can follow every step without opening the repository code.

\section{Setup}

The unit 2-sphere has metric (coordinates $x^1=\theta$, $x^2=\phi$)
\begin{equation}
g = \mathrm{diag}\!\left(g_{\theta\theta},\ g_{\phi\phi}\right)
  = \mathrm{diag}\!\left(1,\ \sin^2\theta\right),
\qquad g^{\theta\theta}=1,\quad g^{\phi\phi}=\frac{1}{\sin^2\theta}.
\end{equation}

\section{Christoffel symbols}

\begin{equation}
\Gamma^\lambda_{\mu\nu} = \tfrac12 g^{\lambda\sigma}\left(
\partial_\mu g_{\sigma\nu} + \partial_\nu g_{\sigma\mu} - \partial_\sigma g_{\mu\nu}\right).
\end{equation}

The only non-vanishing partial derivative is
$\partial_\theta g_{\phi\phi} = 2\sin\theta\cos\theta$. Therefore
\begin{align}
\Gamma^\theta_{\phi\phi} &= -\tfrac12\, g^{\theta\theta}\, \partial_\theta g_{\phi\phi}
= -\sin\theta\cos\theta, \\
\Gamma^\phi_{\theta\phi} = \Gamma^\phi_{\phi\theta}
&= \tfrac12\, g^{\phi\phi}\, \partial_\theta g_{\phi\phi}
= \frac{\cos\theta}{\sin\theta} = \cot\theta.
\end{align}
At $\theta=\pi/4$: $\Gamma^\theta_{\phi\phi}=-1/2$, $\Gamma^\phi_{\theta\phi}=1$.

\section{Riemann curvature}

\begin{equation}
R^\rho_{\ \sigma\mu\nu} = \partial_\mu\Gamma^\rho_{\nu\sigma}
- \partial_\nu\Gamma^\rho_{\mu\sigma}
+ \Gamma^\rho_{\mu\lambda}\Gamma^\lambda_{\nu\sigma}
- \Gamma^\rho_{\nu\lambda}\Gamma^\lambda_{\mu\sigma}.
\end{equation}

The single independent component is
\begin{align}
R^\theta_{\ \phi\theta\phi}
&= \partial_\theta\Gamma^\theta_{\phi\phi} - \Gamma^\theta_{\phi\phi}\Gamma^\phi_{\theta\phi} \\
&= -\left(\cos^2\theta-\sin^2\theta\right) + \sin\theta\cos\theta\cdot\cot\theta \\
&= \sin^2\theta.
\end{align}
At $\theta=\pi/4$: $R^\theta_{\ \phi\theta\phi}=1/2$.

\section{Gauss curvature, Ricci, and Einstein}

Gauss curvature:
\begin{equation}
K = \frac{R^\theta_{\ \phi\theta\phi}}{g_{\phi\phi}}
= \frac{\sin^2\theta}{\sin^2\theta} = 1.
\end{equation}

Ricci tensor (in 2D, $R_{\mu\nu}=K g_{\mu\nu}$):
\begin{equation}
R_{\theta\theta}=K g_{\theta\theta}=1,\qquad
R_{\phi\phi}=K g_{\phi\phi}=\sin^2\theta=\tfrac12 \text{ at }\theta=\pi/4.
\end{equation}

Ricci scalar:
\begin{equation}
R = g^{\mu\nu}R_{\mu\nu} = 2K = 2.
\end{equation}

Einstein tensor (vanishes identically in 2D):
\begin{equation}
G_{\mu\nu} = R_{\mu\nu} - \tfrac12 R g_{\mu\nu} = 0.
\end{equation}

All of these are reproduced numerically by \texttt{EinsteinTensorBuilder} to within
$10^{-3}$. A non-trivial $G_{\mu\nu}$ requires dimension $\ge3$; the unit 3-sphere gives
$G_{\mu\nu}=-g_{\mu\nu}$ (Appendix~\ref{app:tests}).

%=============================================================================

\chapter{Research Programs}
\label{app:programs}

The derivation program is organized into research programs, each closed by executable
experiments. The four programs that carry the technical weight of this monograph --- the
gauge, Koide, continuum, and metric programs --- are narrated in
Chapter~\ref{ch:journey}; this appendix records the program ledger.

\begin{center}
\begin{tabular}{llc}
\toprule
Program & Scope & Experiments \\
\midrule
DATA & Cosmology \& RAR & 10 (DATA-001$\to$010) \\
QM & Quantum foundations & 5 (QM-001$\to$005) \\
QG & Quantum gravity & 31 (QG-001$\to$031) \\
XC & Continuum-limit bridge & (audits + tests) \\
\bottomrule
\end{tabular}
\end{center}

Key results of the full program:
\begin{itemize}
\item $G$ is not fundamental: $G=\ell^2c^3/\hbar$ (gravity emerges from the triple).
\item $(\ell,\tau,\hbar)$ is the irreducible triple — one process, three aspects.
\item Oscillation is a logical inevitability (cannot be removed without destroying $Q$).
\item Gravity is a phase-gradient phenomenon; attraction dominates, repulsion is dark
energy.
\item Gravity manipulation is not possible (QG-023$\to$031).
\item RAR derivation $g_\dagger=cH_0/(2\pi)$ with zero free parameters.
\item Parameter compression SM+GR ($\sim$26) $\to$ TQM ($2{+}3$), $\sim$5--8$\times$
reduction.
\item Novel predictions: $w(z)=-1+0.015(1+z)^{3/2}$; evolving $g_\dagger(z)$; Euclid
sensitivity.
\end{itemize}

\chapter{Development Ledger (Selected)}
\label{app:ledger}

This appendix records, in compact form, a selection of the experiments and audits that
underpin the narrative of Chapter~\ref{ch:timeline}. It is not exhaustive: the full
laboratory book lives in the project's documentation. Each entry states the identifier,
the question or claim, and the recorded outcome. Entries are grouped by era.

\begin{longtable}{p{4.2cm}p{5.2cm}p{5.6cm}}
\toprule
Identifier & Question / Claim & Outcome / Disposition \\
\midrule
\endhead
\multicolumn{3}{l}{\emph{Early oscillator experiments}} \\
\midrule
TQM-001 & Does synchronization emerge? & Yes: order parameter $R\to1$. \\
TQM-002 & Do random matrices give emergence? & No: Wigner spectrum, no structured emergence. \\
TQM-003 & Is static topology sufficient? & No: dynamics outweigh static structure. \\
TQM-004 & Does a field synchronize two asymmetric oscillators? & No: field acts as a resonance medium. \\
TQM-005 & Do resonance clusters form? & Yes: two clusters of three oscillators, $\sim5\%$ energy participation. \\
TQM-006 & Is there a critical resonance density? & Yes: $\rho_c\approx0.09$, percolation-like transition. \\
TQM-007 & How many resonance families exist? & Five reported. \\
TQM-008 & Which families are reproducible? & Only F4 (74\%, score 0.721); the rest are seed artifacts. \\
\midrule
\multicolumn{3}{l}{\emph{Alternative foundations (ResearchX)}} \\
\midrule
X001 & Hidden-assumption audit & 11 found; the static graph is \#1. \\
X002--X004 & Do dynamic graphs innovate? & No: node motion, mobility, growth add nothing open-ended. \\
X005--X006 & Does nonlinearity preserve the quantum reading? & No: eigenmodes break into solitons; linear-algebra artifacts. \\
X007--X008 & Universal species theory & 16 carrier classes across five regimes. \\
X009--X015 & Deepest foundations & Reversibility $+$ self-consistency is minimally sufficient; Rev$\cap$SC $=$ QM. \\
X016--X017 & Reality map & $(R,S)$ is a map, not a dynamical flow. \\
X018--X022 & Complexity staircase & Level 6 not observed; operator evolution needed. \\
X023--X027 & Open-ended evolution search & Delayed saturation; pigeonhole proof for finite systems. \\
X028--X029 & Finite-universe consequences & Saturation physically irrelevant; hybrid architectures optimal. \\
X030--X031 & Quantum optimality & QM is the unique global maximum of finite complexity. \\
X032--X034 & Unification & $L_Q$ is not fundamental; emergence gap closed. \\
X035--X036 & Origin of $Q$; complexity theorem & $Q$ irreducible; Schr\"odinger derived (classification C). \\
X037--X039 & Born, measurement, randomness & Born derived ($\alpha=2$); Many-Worlds incompatible; randomness irreducible. \\
X040--X042 & Time, gravity, dimensionality & Time $=$ partial order; gravity $=$ causal set; $3{+}1$ derived. \\
X043--X046 & Fundamental scales & $G=\ell^2c^3/\hbar$; $\hbar=1$; $c,G,\hbar$ unified; $\Lambda\sim H^2$. \\
X047--X049 & Matter and gauge structure & Particles $=$ defects; gauge structure emerges; selection gap. \\
X051--X056, X059--X060 & Particle physics & Generations, masses, mixing, neutrinos; SM group preferred not derived. \\
X057--X058, X060b--X060f & Parameter compression & Two primitives $+$ $M^2$; $M^2$ irreducible. \\
\midrule
\multicolumn{3}{l}{\emph{Late phases (140--159)}} \\
\midrule
Phase 140 & Causal-set $\Lambda$ & $\Lambda\sim1/\sqrt N$ (postdiction). \\
Phase 144 & Radial acceleration relation & $g_\dagger=cH_0/2\pi$; the $2\pi$ is a numerical accident. \\
Phase 148 & Flavor/Yukawa reducibility & Koide $Q=2/3$ contingent, not derived. \\
Phase 149 & Gauge origin & $U(1)$ derived; $SU(2)$ emergent; $SU(3)$ contingent. \\
Phase 150 & Multiplicity ``3'' & Spacetime 3 derived; internal 3 selected. \\
Phase 151 & Upper bound $N\le3$ & Empirical/contingent; no stability/anomaly/representation bound. \\
Phase 152 & Random Actualization ensembles & Four independent ensembles under a log-normal law. \\
Phase 153 & Cascade unification & Three universality classes are independent cascades. \\
Phase 155 & Neutrino-Koide & Falsified ($Q_{\max}=0.585/0.500<2/3$). \\
Phase 159 & Koide closure & Seven routes, six no-gos, one blocked; closed question. \\
QG-080--100 & Causal-set foundational program & $\Lambda\sim H^2/c^2$; $a_0\sim cH$; terminal answer reached. \\
Hostile review (Round~1) & v1.0 paper audit & 1 valid, 11 partially valid, 0 invalid; verdict NOT READY. \\
Hostile review (Round~2) & FATAL issues audit & 6 resolved, 6 partially resolved, 0 open. \\
v1.0 release & Version readiness & Four dispositions; v1.0 released. \\
\bottomrule
\end{longtable}

\chapter{Glossary}
\label{app:glossary}

\begin{center}
\begin{tabular}{lp{10cm}}
\toprule
Term & Meaning \\
\midrule
$Q$ & principle of individuation (ontology) \\
Random Actualization & ontological chance (assumption A-03) \\
$(\ell,\tau,\hbar)$ & spacetime scale, clock, action (units) \\
$M^2$ & nonlinearity parameter (single continuous) \\
$L_Q$ & graph Laplacian $D-A$ \\
$L_W$ & weighted Laplacian $D_K-K$ \\
BDG & Benincasa--Dowker--Glaser d'Alembertian \\
$\Delta_g$ & Laplace--Beltrami operator \\
$G_{\mu\nu}$ & Einstein tensor $R_{\mu\nu}-\tfrac12 R g_{\mu\nu}$ \\
\derived & computable from primitives by theorem \\
\realunder & real structure, underivable origin \\
\drawn & coincidental draw of the abundance law \\
T-08\dots T-12 & conditional no-go theorems \\
A-03, A-06, A-07, A-10 & assumptions (random actualization; $Z_2$ lift; gauge dynamics; $N$/$n$ split) \\
\bottomrule
\end{tabular}
\end{center}

%=============================================================================

\chapter{References}
\label{app:references}

\begin{thebibliography}{99}

\bibitem{malament1977}
D.~B.~Malament, ``The class of continuous timelike curves determines the topology of
spacetime,'' \emph{J. Math. Phys.} \textbf{18}, 1399 (1977).

\bibitem{hkm1976}
S.~W.~Hawking, A.~R.~King, P.~J.~McCarthy, ``A new topology for curved space--time which
incorporates the causal, differential, and conformal structures,'' \emph{J. Math. Phys.}
\textbf{17}, 174 (1976).

\bibitem{myrheim1978}
J.~Myrheim, ``Statistical geometry,'' CERN-TH-2538 (1978).

\bibitem{bdg2010}
D.~M.~T.~Benincasa, F.~Dowker, ``Scalar curvature of a causal set,'' \emph{Phys. Rev. Lett.}
\textbf{104}, 181301 (2010); L.~Glaser, \emph{Class. Quant. Grav.} \textbf{31}, 095007
(2014).

\bibitem{gleason1957}
A.~M.~Gleason, ``Measures on the closed subspaces of a Hilbert space,''
\emph{J. Math. Mech.} \textbf{6}, 885 (1957).

\bibitem{sorkin}
R.~D.~Sorkin, ``Causal sets: discrete gravity,'' \emph{Proceedings of the Valdivia Summer
School} (2003).

\bibitem{belkin2003}
M.~Belkin, P.~Niyogi, ``Laplacian eigenmaps for dimensionality reduction and data
representation,'' \emph{Neural Computation} \textbf{15}, 1373 (2003).

\bibitem{koide1981}
Y.~Koide, ``Fermion mass relation and the generation structure,'' \emph{Phys. Rev. D}
\textbf{28}, 252 (1983) [letter 1981].

\bibitem{bertrand}
J.~Bertrand, ``Th\'eor\`eme relatif au mouvement d'un point attir\'e vers un centre fixe,''
\emph{C. R. Acad. Sci.} \textbf{77}, 849 (1873).

\end{thebibliography}

\end{document}
